\documentclass[11pt,oneside,headings=big]{scrreprt}

% Shared visual style for the repository's LaTeX documents.
% Keep document-specific packages, metadata, and content in each source file.

\usepackage[letterpaper,margin=0.86in,headheight=16pt,footskip=28pt]{geometry}
\usepackage[T1]{fontenc}
\usepackage{newtxtext,newtxmath}
\usepackage{microtype}
\usepackage{graphicx}
\usepackage{booktabs,tabularx,longtable,array}
\usepackage[table]{xcolor}
\usepackage[most]{tcolorbox}
\usepackage{enumitem}
\usepackage{scrlayer-scrpage}
\usepackage{caption}
\usepackage{subcaption}
\usepackage{hyperref}

\definecolor{ImuNavy}{HTML}{17324D}
\definecolor{ImuBlue}{HTML}{2E6F9E}
\definecolor{ImuTeal}{HTML}{168A8A}
\definecolor{ImuOrange}{HTML}{D9792B}
\definecolor{ImuGold}{HTML}{D9A441}
\definecolor{ImuInk}{HTML}{1E252B}
\definecolor{ImuSlate}{HTML}{52606D}
\definecolor{ImuPale}{HTML}{EEF4F7}
\definecolor{ImuPaleTeal}{HTML}{EAF6F5}
\definecolor{ImuPaleGold}{HTML}{FFF7E6}
\definecolor{ImuRule}{HTML}{CBD7DE}
\definecolor{ImuCode}{HTML}{F5F7F8}

\setkomafont{disposition}{\color{ImuNavy}\bfseries}
\setkomafont{chapter}{\color{ImuNavy}\bfseries\LARGE}
\setkomafont{section}{\color{ImuBlue}\bfseries\Large}
\setkomafont{subsection}{\color{ImuTeal}\bfseries\large}
\setkomafont{caption}{\small}
\setkomafont{captionlabel}{\color{ImuNavy}\bfseries}
\renewcommand*{\chapterformat}{%
    \colorbox{ImuNavy}{\parbox[c][1.65em][c]{2.35em}{\centering\color{white}\thechapter}}\enskip
}

\clearpairofpagestyles
\ihead{\textsc{IMU Error Model}}
\ohead{\headmark}
\cfoot*{\pagemark}
\addtokomafont{pageheadfoot}{\small\color{ImuSlate}}
\automark[section]{chapter}

\setlist[itemize]{leftmargin=1.4em,itemsep=0.25em,topsep=0.35em}
\setlist[enumerate]{leftmargin=1.6em,itemsep=0.25em,topsep=0.35em}
\renewcommand{\arraystretch}{1.18}
\setlength{\tabcolsep}{3pt}
\emergencystretch=1.5em

\newcolumntype{Y}{>{\raggedright\arraybackslash}X}
\newcolumntype{L}[1]{>{\raggedright\arraybackslash}p{#1}}
\newcolumntype{C}[1]{>{\centering\arraybackslash}p{#1}}
\newcommand{\software}[1]{\texttt{#1}}

\usepackage{mathtools,bm}
\usepackage{siunitx}
\usepackage{listings}
\usepackage[ruled,vlined,linesnumbered]{algorithm2e}
\usepackage{tikz}
\usetikzlibrary{arrows.meta,positioning,calc,fit,backgrounds,shapes.geometric,matrix}
\usepackage{pgfplots}
\pgfplotsset{compat=1.18}
\usepackage{csquotes}
\usepackage[backend=biber,style=ieee,sorting=none,maxbibnames=8]{biblatex}
\addbibresource{imu_error_model_references.bib}
\usepackage[nameinlink,noabbrev]{cleveref}
\usepackage{bookmark}

\hypersetup{
    colorlinks=true,
    linkcolor=ImuBlue,
    citecolor=ImuTeal,
    urlcolor=ImuOrange,
    pdftitle={IMU Error Model: Mathematical Specification, Software Contract, and Verification Guide},
    pdfauthor={imu-error-model project},
    pdfsubject={Three-axis accelerometer and gyroscope error simulation},
    pdfkeywords={IMU, accelerometer, gyroscope, Allan deviation, Gauss-Markov, sensor error}
}

\sisetup{
    detect-all,
    per-mode=symbol,
    exponent-product=\mathord{\times},
    range-phrase=--,
    range-units=single
}
\DeclareSIUnit\gzero{g_0}
\DeclareSIUnit\ppm{ppm}
\DeclareSIUnit\arcsec{arcsec}

\newtcolorbox{boundarybox}[1][]{
    enhanced,breakable,
    colback=ImuPaleGold,colframe=ImuGold,
    boxrule=0.8pt,arc=2pt,
    left=8pt,right=8pt,top=7pt,bottom=7pt,
    title={Model boundary},
    fonttitle=\bfseries\color{ImuNavy},
    #1
}
\newtcolorbox{implementationbox}[1][]{
    enhanced,breakable,
    colback=ImuPaleTeal,colframe=ImuTeal,
    boxrule=0.8pt,arc=2pt,
    left=8pt,right=8pt,top=7pt,bottom=7pt,
    title={Normative implementation behavior},
    fonttitle=\bfseries\color{ImuNavy},
    #1
}
\newtcolorbox{engineeringbox}[1][]{
    enhanced,breakable,
    colback=ImuPale,colframe=ImuBlue,
    boxrule=0.8pt,arc=2pt,
    left=8pt,right=8pt,top=7pt,bottom=7pt,
    title={Engineering note},
    fonttitle=\bfseries\color{ImuNavy},
    #1
}
\newtcolorbox{warningbox}[1][]{
    enhanced,breakable,
    colback=white,colframe=ImuOrange,
    boxrule=1.0pt,arc=2pt,
    borderline west={3pt}{0pt}{ImuOrange},
    left=9pt,right=8pt,top=7pt,bottom=7pt,
    title={Interpretation warning},
    fonttitle=\bfseries\color{ImuOrange},
    #1
}

\lstdefinestyle{imuPython}{
    language=Python,
    backgroundcolor=\color{ImuCode},
    basicstyle=\ttfamily\small,
    keywordstyle=\color{ImuBlue}\bfseries,
    commentstyle=\color{ImuSlate}\itshape,
    stringstyle=\color{ImuTeal},
    numbers=left,
    numberstyle=\tiny\color{ImuSlate},
    numbersep=8pt,
    frame=single,
    rulecolor=\color{ImuRule},
    framerule=0.5pt,
    breaklines=true,
    showstringspaces=false,
    tabsize=4,
    columns=fullflexible
}

\newcommand{\Rwb}[1][]{\mathbf{R}^{w}_{b#1}}
\newcommand{\Rbw}[1][]{\mathbf{R}^{b#1}_{w}}
\newcommand{\Exp}{\operatorname{Exp}}
\newcommand{\Log}{\operatorname{Log}}
\newcommand{\veeop}{(\cdot)^\vee}
\newcommand{\skewmat}[1]{\left[#1\right]_{\times}}
\newcommand{\Normal}{\mathcal{N}}
\newcommand{\E}{\mathbb{E}}
\newcommand{\Cov}{\operatorname{Cov}}
\newcommand{\Var}{\operatorname{Var}}
\newcommand{\diag}{\operatorname{diag}}
\newcommand{\tr}{\operatorname{tr}}
\newcommand{\clip}{\operatorname{clip}}
\newcommand{\round}{\operatorname{round}}
\newcommand{\sat}{\operatorname{sat}}
\newcommand{\Had}{\odot}
\newcommand{\dd}{\,\mathrm{d}}
\newcommand{\transpose}{\mathsf{T}}
\newcommand{\vect}[1]{\bm{#1}}
\newcommand{\mat}[1]{\bm{#1}}
\newcommand{\field}[1]{\texttt{#1}}
\newcommand{\pkg}{\texttt{imu-error-model}}

\title{IMU Error Model}
\subtitle{Mathematical Specification, Software Contract, and Verification Guide}
\author{\pkg{} project}
\date{\today}

\begin{document}
    \hypersetup{pageanchor=false}

    \begin{titlepage}
        \begin{tikzpicture}[remember picture,overlay]
    \fill[ImuNavy] (current page.north west) rectangle ($(current page.north east)+(0,-1.35in)$);
    \fill[ImuTeal] ($(current page.south west)+(0,0.34in)$) rectangle (current page.south east);
    \draw[ImuRule,line width=0.8pt] ($(current page.north west)+(0.86in,-2.02in)$) -- ($(current page.north east)+(-0.86in,-2.02in)$);

    \coordinate (O) at ($(current page.north east)+(-1.75in,-0.77in)$);
    \draw[-{Latex[length=3.2mm]},line width=1.6pt,white] (O) -- ++(0.72,0) node[right,font=\small\bfseries] {$x_b$};
    \draw[-{Latex[length=3.2mm]},line width=1.6pt,ImuGold] (O) -- ++(-0.30,0.58) node[above left,font=\small\bfseries] {$y_b$};
    \draw[-{Latex[length=3.2mm]},line width=1.6pt,ImuOrange] (O) -- ++(0,-0.67) node[below,font=\small\bfseries] {$z_b$};
        \end{tikzpicture}

        \vspace*{1.58in}
        {\color{ImuNavy}\fontsize{31}{35}\selectfont\bfseries IMU Error Model\par}
        \vspace{0.28cm}
        {\color{ImuSlate}\fontsize{17}{22}\selectfont Mathematical Specification, Software Contract,\par
        and Verification Guide\par}

        \vspace{0.78in}
        \begin{tcolorbox}[
            width=\textwidth,
            colback=ImuPale,
            colframe=ImuRule,
            boxrule=0.6pt,
            arc=2pt,
            left=12pt,right=12pt,top=11pt,bottom=11pt]
            \large
            A frame-explicit, reproducible model for deterministic and stochastic errors in
            three-axis accelerometer and gyroscope increments.
        \end{tcolorbox}

    \end{titlepage}
    \clearpage
    \hypersetup{pageanchor=true}
    \pagenumbering{roman}


    \chapter*{Abstract}
    \addcontentsline{toc}{chapter}{Abstract}
    The \pkg{} package converts an externally generated truth trajectory into
    accelerometer and gyroscope increments corrupted by configurable deterministic
    and stochastic sensor effects. The model is intentionally narrower than a
    rigid-body simulator or inertial navigation system: it neither advances vehicle
    dynamics nor reconstructs position, velocity, or attitude. Its responsibility is
    to preserve a precise frame and sampling contract while applying scale-factor
    error, nonlinearity, thermal effects, run-level misalignment, turn-on bias,
    stationary Gauss--Markov drift, random walk, finite-band flicker-like bias,
    correlated white noise, saturation, quantization, and output scaling.

    This specification formalizes that transformation on the rotation group
    $\mathrm{SO}(3)$, derives exact discrete-time stochastic updates, records the
    units and state lifecycle, and connects model parameters to Allan-deviation and
    power-spectral interpretations. It also supplies sensitivity Jacobians,
    downstream navigation-error approximations, profile-construction guidance, and
    a verification strategy suitable for deterministic unit tests and statistical
    Monte Carlo regression. The result is both a mathematical reference for model
    review and a practical contract for simulation, estimator, and controller
    integration.

    \begin{boundarybox}
        The package corrupts truth measurements. It does not create truth, add gravity,
        propagate an INS, estimate states, or certify a hardware device. Hardware-oriented
        profiles included in this repository are notional research approximations,
        estimated from public datasheets or other cited source material. Their source
        links identify evidence used for parameter selection; they do not make the
        resulting values authoritative, vendor-approved, certified, or guaranteed.
    \end{boundarybox}

    \tableofcontents
    \listoffigures
    \listoftables


    \chapter*{Notation and conventions}
    \addcontentsline{toc}{chapter}{Notation and conventions}
    \begin{longtable}{@{}L{1.35in}L{1.45in}L{3.35in}@{}}
        \caption{Principal notation used throughout the specification.}\label{tab:notation}\\
        \toprule
        \textbf{Symbol} & \textbf{Type / units} & \textbf{Meaning} \\
        \midrule
        \endfirsthead
        \toprule
        \textbf{Symbol} & \textbf{Type / units} & \textbf{Meaning} \\
        \midrule
        \endhead
        $w,b$ & frames & World frame and body/sensor reference frame. \\
        $\Rwb(t)$ & $\mathrm{SO}(3)$ & Rotation mapping body coordinates at time $t$ into world coordinates. \\
        $\vect v^w(t)$ & \si{\meter\per\second} & Truth velocity supplied by the caller, with gravity excluded by contract. \\
        $\Delta t_k$ & \si{\second} & Sample interval $t_{k+1}-t_k>0$. \\
        $\Delta\vect v_{b,k}$ & \si{\meter\per\second} & Ideal body-start-frame velocity increment. \\
        $\Delta\vect\theta_{b,k}$ & \si{\radian} & Ideal body-start-frame rotation vector. \\
        $\vect u_{a,k}$ & \si{\meter\per\second\squared} & Ideal accelerometer rate $\Delta\vect v_{b,k}/\Delta t_k$. \\
        $\vect u_{g,k}$ & \si{\radian\per\second} & Ideal gyroscope rate $\Delta\vect\theta_{b,k}/\Delta t_k$. \\
        $s\in\{a,g\}$ & index & Accelerometer or gyroscope channel. \\
        $\Had$ & operator & Hadamard, or elementwise, product. \\
        $\skewmat{\vect x}$ & $\mathfrak{so}(3)$ & Skew-symmetric matrix satisfying $\skewmat{\vect x}\vect y=\vect x\times\vect y$. \\
        $\Exp,\Log$ & maps & Exponential and logarithm maps between $\mathfrak{so}(3)$ and $\mathrm{SO}(3)$. \\
        $\vect b_{\mathrm{on},s}$ & channel rate units & Turn-on bias sampled at reset and fixed for one run. \\
        $\vect b_{s,k}$ & channel rate units & In-run Gauss--Markov or random-walk bias state. \\
        $\vect f_{s,k}$ & channel rate units & Finite-band flicker-like bias approximation. \\
        $\vect n_{s,k}$ & channel rate units & White rate-noise sample. \\
        $\mat C_s$ & rate$^2\!\cdot$\si{\second} & Continuous-time three-axis white-noise covariance. \\
        $N_s$ & rate$/\sqrt{\si{\hertz}}$ & Independent-axis white-noise density when $\mat C_s$ is not supplied. \\
        $T$ & \si{\second} & Temperature-independent correlation time when used in a stochastic process; temperature is written $T_{\mathrm C}$. \\
        $\vect\delta_s$ & channel rate units & Three-axis quantization step; a scalar is broadcast. \\
        $L_s$ & channel rate units & Symmetric measurement range. \\
        \bottomrule
    \end{longtable}

    \clearpage
    \pagenumbering{arabic}


    \chapter{Project overview}
    \label{chap:overview}


    \section{Purpose}
    The model provides a controlled boundary between a truth-producing simulation and
    algorithms that should operate on imperfect inertial measurements. A typical
    closed-loop integration is
    \begin{equation}
        \text{equations of motion}
        \longrightarrow \text{truth snapshot}
        \longrightarrow \text{IMU error model}
        \longrightarrow \text{estimator or controller}
        \longrightarrow \text{command},
    \end{equation}
    where the controller never receives an uncorrupted state unless the test
    explicitly asks for a truth-fed baseline.

    The package is designed for robotics, aerospace simulation, inertial-navigation
    research, estimator verification, and sensor-model validation. It is especially
    useful when the equations of motion already exist and the missing capability is a
    repeatable sensor boundary that can expose algorithms to realistic error
    structure rather than a single ad hoc Gaussian perturbation.


    \section{Model boundary}
    \begin{table}[htbp]
        \centering
        \caption{Included and excluded responsibilities.}
        \label{tab:scope}
        \begin{tabularx}{\textwidth}{@{}Y Y@{}}
\toprule
\textbf{In scope} & \textbf{Out of scope} \\
\midrule
Body-frame accelerometer and gyroscope increments & Rigid-body force and moment integration \\
Scale factor, radial nonlinearity, and thermal terms & Gravity, Earth rotation, Coriolis, and transport-rate models \\
Run-level misalignment and turn-on repeatability & Strapdown navigation mechanization or sensor fusion \\
White, correlated, Gauss--Markov, random-walk, and finite-band flicker-like errors & Hardware drivers, timestamp acquisition, or device communication \\
Saturation, quantization, and output scaling & Vendor certification or guaranteed datasheet equivalence \\
Reproducible stochastic state and analysis fixtures & Automatic calibration from laboratory data \\
Protocol-based alternate model implementations & Magnetometer, GNSS, camera, or seeker measurement models \\
\bottomrule
        \end{tabularx}
    \end{table}

    \begin{implementationbox}
        The public model consumes a timestamp, world-frame truth velocity with gravity
        excluded, a proper body-to-world rotation matrix, and an optional temperature.
        The first call initializes history and returns zero increments. Every later call
        requires a strictly increasing timestamp and returns increments associated with
        that exact interval.
    \end{implementationbox}


    \section{Signal path}
    \Cref{fig:pipeline} summarizes the implemented transformation. Deterministic
    terms operate on an ideal interval-average rate, stochastic terms maintain
    state across samples, and output-electronics effects act before conversion back
    to increments.

    \begin{figure}[htbp]
        \centering
        \begin{tikzpicture}[
            node distance=4mm,
            block/.style={
                draw=ImuRule,rounded corners=2pt,minimum height=28mm,
                text width=26mm,align=center,fill=white,
                font=\small,inner sep=5pt
            },
            arrow/.style={-{Latex[length=2.5mm]},line width=0.85pt,draw=ImuBlue}
        ]
  \node[block,fill=ImuPale] (truth) {\textbf{Truth snapshots}\\[2pt]
      $t_k,\,\vect v^w_k,$\\$\Rwb(t_k),\,T_k$};
  \node[block,right=of truth,fill=ImuPale] (kin) {\textbf{Kinematics}\\[2pt]
      $\Delta\vect v_b$\\$\Delta\vect\theta_b$\\divide by $\Delta t$};
  \node[block,right=of kin,fill=ImuPaleTeal] (det) {\textbf{Deterministic channel}\\[2pt]
  thermal bias\\scale and nonlinearity\\run misalignment};
  \node[block,right=of det,fill=ImuPaleGold] (stoch) {\textbf{Stochastic state}\\[2pt]
  turn-on bias\\GM / random walk\\flicker + white noise};
  \node[block,right=of stoch,fill=ImuPale] (out) {\textbf{Output electronics}\\[2pt]
  clip\\quantize\\scale and integrate};
  \draw[arrow] (truth) -- (kin);
  \draw[arrow] (kin) -- (det);
  \draw[arrow] (det) -- (stoch);
  \draw[arrow] (stoch) -- (out);
        \end{tikzpicture}
        \caption{High-level measurement pipeline. The accelerometer and gyroscope use the
        same channel structure but different physical units and independent
        configurations.}
        \label{fig:pipeline}
    \end{figure}


    \section{Design principles}
    The implementation follows five principles.
    \begin{enumerate}
        \item \textbf{Frames are explicit.} Every vector is labeled by the frame in
        which its coordinates are expressed. Rotation direction is encoded in the
        public argument name \field{orientation\_world\_from\_body}.
        \item \textbf{Continuous-time parameters have exact discrete-time mappings.}
        White-noise density and Gauss--Markov correlation time remain meaningful when
        the sample interval changes.
        \item \textbf{Run-level and in-run uncertainty are separate.} Turn-on bias and
        misalignment describe repeatability between resets; drift processes describe
        evolution within one run.
        \item \textbf{Nonlinear output effects are modeled where they occur.}
        Saturation and quantization act on rates before the result is integrated into
        an increment.
        \item \textbf{Statistical behavior is testable.} The package exposes enough
        structure to verify covariance, autocorrelation, Allan deviation, frame
        invariance, and reset reproducibility independently.
    \end{enumerate}


    \chapter{Frames, rotations, and the truth-state contract}
    \label{chap:truth}


    \section{Rotation convention}
    The rotation matrix
    \begin{equation}
        \Rwb(t)\in\mathrm{SO}(3),
        \qquad
        \mathrm{SO}(3)=\left\{\mat R\in\mathbb{R}^{3\times3}
                           \mid \mat R^\transpose\mat R=\mat I,\ \det\mat R=1\right\},
    \end{equation}
    maps body-frame coordinates into world-frame coordinates:
    \begin{equation}
        \vect x^w = \Rwb(t)\vect x^b,
        \qquad
        \vect x^b = \Rwb(t)^\transpose\vect x^w.
        \label{eq:frame-map}
    \end{equation}
    The transpose is the inverse because a proper rotation is orthogonal. This
    convention agrees with the public name
    \field{orientation\_world\_from\_body}; reversing it silently would change the
    sign and frame of both channels.

    For \(\vect\phi=[\phi_x,\phi_y,\phi_z]^\transpose\), define the hat operator
    \begin{equation}
        \skewmat{\vect\phi}=
        \begin{bmatrix}
            0 & -\phi_z & \phi_y\\
            \phi_z & 0 & -\phi_x\\
            -\phi_y & \phi_x & 0
        \end{bmatrix},
        \qquad
        \skewmat{\vect\phi}\vect x=\vect\phi\times\vect x.
    \end{equation}
    Its inverse on skew-symmetric matrices is the vee operator. The exponential map
    is given by Rodrigues' formula
    \begin{equation}
        \Exp\!\left(\skewmat{\vect\phi}\right)
        =\mat I
        +\frac{\sin\theta}{\theta}\skewmat{\vect\phi}
        +\frac{1-\cos\theta}{\theta^2}\skewmat{\vect\phi}^{2},
        \qquad \theta=\lVert\vect\phi\rVert_2,
        \label{eq:rodrigues}
    \end{equation}
    with the limiting series used near \(\theta=0\). Conversely, for a relative
    rotation \(\mat R\),
    \begin{align}
        \theta &= \cos^{-1}\!\left(\frac{\tr(\mat R)-1}{2}\right),\\
        \Log(\mat R)^\vee
        &=\frac{\theta}{2\sin\theta}
        \begin{bmatrix}
            R_{32}-R_{23}\\
            R_{13}-R_{31}\\
            R_{21}-R_{12}
        \end{bmatrix},
        \label{eq:so3-log}
    \end{align}
    away from the numerically delicate neighborhoods of \(0\) and \(\pi\). Robust
    implementations use a small-angle series near zero and a diagonal-based branch
    near \(\pi\) rather than evaluating \cref{eq:so3-log} literally
    ~\cite{sola2017,barfoot2017}.


    \section{Ideal interval increments}
    Let the interval begin at \(t_k\) and end at \(t_{k+1}\), with
    \begin{equation}
        \Delta t_k=t_{k+1}-t_k>0.
    \end{equation}
    The ideal increments are
    \begin{align}
        \Delta\vect v_{b,k}
        &=\Rwb(t_k)^\transpose
        \left(\vect v^w(t_{k+1})-\vect v^w(t_k)\right),
        \label{eq:ideal-dv}\\
        \Delta\vect\theta_{b,k}
        &=\Log\!\left(\Rwb(t_k)^\transpose\Rwb(t_{k+1})\right)^\vee.
        \label{eq:ideal-dtheta}
    \end{align}
    Both are expressed in the body axes at the \emph{start} of the interval. The
    corresponding interval-average rates are
    \begin{equation}
        \vect u_{a,k}=\frac{\Delta\vect v_{b,k}}{\Delta t_k},
        \qquad
        \vect u_{g,k}=\frac{\Delta\vect\theta_{b,k}}{\Delta t_k}.
        \label{eq:ideal-rates}
    \end{equation}

    \begin{warningbox}
        \(\Delta\vect v_{b,k}\) is not a world-frame velocity difference, and
        \(\Delta\vect\theta_{b,k}\) is not a world-frame rotation vector. To express
        an ideal or measured start-frame increment in world axes, multiply it by
        \(\Rwb(t_k)\).
    \end{warningbox}

    The rotation increment in \cref{eq:ideal-dtheta} is coordinate-invariant with
    respect to the world frame: premultiplying every attitude by one fixed world
    rotation \(\mat S\in\mathrm{SO}(3)\) leaves
    \begin{equation}
        (\mat S\Rwb(t_k))^\transpose(\mat S\Rwb(t_{k+1}))
        =\Rwb(t_k)^\transpose\Rwb(t_{k+1})
    \end{equation}
    unchanged. This identity is a high-value frame regression test.


    \section{Gravity-excluded velocity contract}
    The package does not evaluate a gravity model. The caller supplies
    \(\vect v^w\) under the documented convention that gravity has already been
    excluded from the truth quantity being differenced. Consequently, the model does
    not add or subtract \(\vect g\), and it does not claim to reproduce a complete
    specific-force mechanization.

    \begin{boundarybox}
        The input contract is intentionally simulation-facing. A physical accelerometer
        measures specific force, while this package receives an already prepared truth
        velocity signal and corrupts its interval change. The equations of motion or
        truth adapter is responsible for making that signal consistent with the intended
        experiment.
    \end{boundarybox}


    \section{Sampling semantics}
    The first \software{measure} call establishes the previous timestamp, velocity,
    and orientation and returns zero increments. A later call is valid only when
    \(t_{k+1}>t_k\). The output records
    \field{start\_time}, \field{end\_time}, and the derived \field{dt}; the reported
    \field{delta\_v} and \field{delta\_theta} correspond exactly to that interval.

    Endpoint differencing makes the contract easy to integrate with a time-stepped
    simulation, but it does not perform explicit coning or sculling compensation.
    When the body rotates substantially inside one interval, the caller should sample
    at a sufficiently high rate or supply a specialized model that has access to
    substep truth.


    \section{Input validity}
    A conforming implementation should reject or clearly diagnose
    \begin{equation}
        \Delta t_k\le 0,
        \qquad
        \lVert\mat R^\transpose\mat R-\mat I\rVert_F>\varepsilon_R,
        \qquad
        \det\mat R\le 0,
    \end{equation}
    as well as non-finite timestamps, vectors, matrices, temperatures, or
    configuration values. If small orthogonality drift is tolerated, any projection
    back to \(\mathrm{SO}(3)\) must be explicit; silently repairing arbitrary input
    can conceal upstream simulation faults.

    \begin{table}[htbp]
        \centering
        \caption{Public sample quantities and canonical units.}
        \label{tab:public-units}
        \begin{tabularx}{\textwidth}{@{}L{1.9in}L{1.7in}Y@{}}
\toprule
\textbf{Quantity} & \textbf{Unit} & \textbf{Interpretation} \\
\midrule
Timestamp & \si{\second} & Monotonic simulation time. \\
Truth velocity & \si{\meter\per\second} & World-frame, gravity-excluded input. \\
Orientation & dimensionless & Proper body-to-world rotation matrix. \\
Temperature & \si{\degreeCelsius} & Optional per-sample environmental input. \\
\field{delta\_v} & \si{\meter\per\second} & Corrupted velocity increment in body axes at interval start. \\
\field{delta\_theta} & \si{\radian} & Corrupted rotation vector in body axes at interval start. \\
\bottomrule
        \end{tabularx}
    \end{table}


    \chapter{Implemented measurement model}
    \label{chap:measurement}


    \section{Channel-generic notation}
    The accelerometer and gyroscope share one mathematical pipeline. Let
    \(s\in\{a,g\}\) identify the channel, with ideal rate
    \begin{equation}
        \vect u_{s,k}=\begin{cases}
                          \vect u_{a,k},&s=a,\\
                          \vect u_{g,k},&s=g.
        \end{cases}
    \end{equation}
    All vectors in this chapter are three-axis sensor quantities. Multiplication is
    ordinary matrix multiplication unless \(\Had\) is shown.

    Several configuration fields accept either one scalar or an explicit three-axis
    value. A scalar is broadcast to all axes. In the current schema, scale-factor
    error, turn-on bias standard deviation, in-run bias standard deviation, and
    quantization step support this form. Anisotropic reset misalignment is expressed
    by a positive-semidefinite rotation-vector covariance; the scalar
    \field{misalignment\_std} remains an isotropic shortcut. White-noise
    anisotropy and cross-axis correlation use \field{noise\_covariance}, while
    thermal coefficients already support scalar or three-axis values.


    \section{Thermal state}
    For channel \(s\), define the temperature offset
    \begin{equation}
        \Delta T_{s,k}=T_k-T_{0,s},
    \end{equation}
    and the three configured thermal responses
    \begin{align}
        \vect\beta_{T,s,k}&=\vect c_{b,s}\,\Delta T_{s,k},
        &&\text{thermal bias offset},\label{eq:thermal-bias}\\
        \vect\lambda_{T,s,k}&=\vect 1+\vect c_{n,s}\,\Delta T_{s,k},
        &&\text{white-noise multiplier},\label{eq:thermal-noise}\\
        \vect\kappa_{T,s,k}&=\vect c_{\mathrm{sf},s}\,\Delta T_{s,k},
        &&\text{scale-factor offset}.\label{eq:thermal-scale}
    \end{align}
    A scalar coefficient is broadcast to all axes; a three-vector acts elementwise.
    The coefficient units are channel-rate units per \si{\degreeCelsius} for
    \(\vect c_b\), and \si{\per\degreeCelsius} for \(\vect c_n\) and
    \(\vect c_{\mathrm{sf}}\). The default reference temperature is
    \SI{25}{\degreeCelsius}.

    If temperature is unavailable, the thermal model returns the identity state
    \begin{equation}
        \vect\beta_{T,s,k}=\vect 0,
        \qquad
        \vect\lambda_{T,s,k}=\vect 1,
        \qquad
        \vect\kappa_{T,s,k}=\vect 0.
        \label{eq:thermal-identity}
    \end{equation}
    The default linear model requires \(\vect\lambda_{T,s,k}\ge\vect0\)
    componentwise so that a nominal standard deviation is not multiplied by a
    negative factor.


    \section{Scale factor and radial nonlinearity}
    The implemented deterministic channel map is
    \begin{equation}
        \vect d_s(\vect u_{s,k},T_k)=
        \vect u_{s,k}\Had
        \left(\vect 1+\vect s_s+\vect\kappa_{T,s,k}\right)
        \Had\left(\vect 1+n_s\lVert\vect u_{s,k}\rVert_2^2\vect 1\right)
        +\vect\beta_{T,s,k}.
        \label{eq:deterministic-map}
    \end{equation}
    Here \( \vect s_s \) is the fixed three-axis scale-factor coefficient and
    \(n_s\) is a radial nonlinearity coefficient. A scalar scale factor is
    broadcast to all axes. The radial form preserves the direction-dependent axis
    scaling while making the nonlinear gain depend on the total input magnitude.

    It is useful to define
    \begin{equation}
        \mat A_{s,k}=\diag\!\left(\vect1+\vect s_s+\vect\kappa_{T,s,k}\right),
        \qquad
        h_s(\vect u)=1+n_s\vect u^\transpose\vect u,
    \end{equation}
    so that
    \begin{equation}
        \vect d_s=\mat A_{s,k}\vect u_{s,k}h_s(\vect u_{s,k})
        +\vect\beta_{T,s,k}.
        \label{eq:deterministic-matrix}
    \end{equation}
    The local input Jacobian is then
    \begin{equation}
        \frac{\partial\vect d_s}{\partial\vect u}
        =\mat A_{s,k}
        \left[
            h_s(\vect u)\mat I+2n_s\vect u\vect u^\transpose
        \right].
        \label{eq:deterministic-jacobian}
    \end{equation}
    This derivative is useful for calibration, sensitivity analysis, and explaining
    why a large input can amplify even a modest nonlinear coefficient.


    \section{Run-level misalignment}
    At reset, the channel samples a small rotation vector
    \begin{equation}
        \vect\epsilon_s\sim\Normal(\vect0,\mat\Sigma_{\epsilon,s}),
        \qquad
        \mat\Sigma_{\epsilon,s}=m_s^2\mat I
        \text{ for the isotropic scalar shortcut.}
    \end{equation}
    and constructs
    \begin{equation}
        \mat M_s=\Exp\!\left(\skewmat{\vect\epsilon_s}\right).
        \label{eq:misalignment}
    \end{equation}
    The matrix remains fixed for the run. For small angles,
    \begin{equation}
        \mat M_s\vect d_s
        \approx \vect d_s+\vect\epsilon_s\times\vect d_s
        =\vect d_s-\skewmat{\vect d_s}\vect\epsilon_s,
    \end{equation}
    so the first-order sensitivity to misalignment is
    \begin{equation}
        \left.\frac{\partial(\mat M_s\vect d_s)}{\partial\vect\epsilon_s}
        \right|_{\vect\epsilon_s=\vect0}
        =-\skewmat{\vect d_s}.
        \label{eq:misalignment-jacobian}
    \end{equation}
    A rotation preserves vector norm but couples signal from one axis into the
    others; this is physically different from independent scale-factor error.


    \section{Additive stochastic rate error}
    After deterministic distortion and run-level misalignment, the model adds
    \begin{equation}
        \vect e_{s,k}
        =\vect b_{\mathrm{on},s}
        +\vect b_{s,k}
        +\vect f_{s,k}
        +\vect n_{s,k},
        \label{eq:additive-error}
    \end{equation}
    which gives the pre-output rate
    \begin{equation}
        \vect r_{s,k}
        =\mat M_s\vect d_s(\vect u_{s,k},T_k)+\vect e_{s,k}.
        \label{eq:pre-output-rate}
    \end{equation}
    The stochastic terms are defined in \cref{chap:stochastic}. They are sampled in
    sensor axes and are not rotated into the world frame.


    \section{Saturation and quantization}
    For a symmetric rate range \(L_s\),
    \begin{equation}
        \bar{\vect r}_{s,k}
        =\clip(\vect r_{s,k},-L_s,L_s)
        \label{eq:clip}
    \end{equation}
    componentwise. If no range is configured or \field{apply\_clipping} is false,
    clipping is the identity. With a
    positive three-axis quantization step \(\vect\delta_s\), where a
    scalar value is broadcast to all axes,
    \begin{equation}
        \vect q_{s,k}
        =\vect\delta_s\Had\round\!\left(\frac{\bar{\vect r}_{s,k}}{\vect\delta_s}\right),
        \label{eq:quantization}
    \end{equation}
    and otherwise \(\vect q_{s,k}=\bar{\vect r}_{s,k}\).

    Away from saturation and assuming round-to-nearest, quantization introduces the
    componentwise bound
    \begin{equation}
        \lvert q_{s,k,i}-r_{s,k,i}\rvert\le\frac{\delta_{s,i}}{2}.
        \label{eq:quant-bound}
    \end{equation}
    After integration, the corresponding increment bound is
    \begin{equation}
        \left\lvert\Delta y^{\mathrm{quant}}_{s,k,i}\right\rvert
        \le \Delta t_k\,\lvert o_sO_s\rvert\frac{\delta_{s,i}}{2}.
    \end{equation}
    The exact quantizer is piecewise constant and therefore not differentiable at
    thresholds; any smooth surrogate used in optimization is an approximation, not
    the physical output rule.


    \section{Reported increments}
    Let \(o_s\) denote the channel output scale and \(O_s\) the corresponding global
    \software{ImuConfig} output scale. The reported increments are
    \begin{align}
        \Delta\vect v^{\mathrm{meas}}_k
        &=\Delta t_k\,o_aO_a\vect q_{a,k},\label{eq:reported-dv}\\
        \Delta\vect\theta^{\mathrm{meas}}_k
        &=\Delta t_k\,o_gO_g\vect q_{g,k}.
        \label{eq:reported-dtheta}
    \end{align}
    Combining the stages gives the compact channel operator
    \begin{equation}
        \boxed{
            \Delta\vect y^{\mathrm{meas}}_{s,k}
            =\Delta t_k\,o_sO_s\,
            \mathcal{Q}_{\vect\delta_s}\!\left(
                                             \mathcal{C}_{L_s}\!\left[
                                                                    \mat M_s\vect d_s(\vect u_{s,k},T_k)
                    +\vect b_{\mathrm{on},s}+\vect b_{s,k}+\vect f_{s,k}+\vect n_{s,k}
                \right]\right)
        }
        \label{eq:unified-model}
    \end{equation}
    where \(\mathcal{C}\) and \(\mathcal{Q}\) are the componentwise clip and
    quantization operators. Equation~\eqref{eq:unified-model} is the central
    normative model.

    \begin{implementationbox}
        The final output is an increment obtained from a corrupted rate. It is generally
        not equal to the ideal kinematic increment plus one independently sampled
        increment error because scale, nonlinearity, misalignment, clipping, and
        quantization are nonlinear or multiplicative operations.
    \end{implementationbox}


    \section{Thermal sensitivity}
    For the linear thermal model and fixed \(\vect u\), the deterministic rate
    sensitivity is
    \begin{equation}
        \frac{\partial\vect d_s}{\partial T}
        =h_s(\vect u)\diag(\vect u)\vect c_{\mathrm{sf},s}
        +\vect c_{b,s}.
        \label{eq:thermal-jacobian}
    \end{equation}
    Temperature also changes the white-noise covariance through
    \(\vect\lambda_{T,s,k}\). Therefore a thermal sweep can alter both the mean and
    the scatter of the measured output; validating only the average response is
    insufficient.


    \chapter{Stochastic error processes}
    \label{chap:stochastic}


    \section{Separation by time scale}
    A useful IMU model separates errors by how they behave in time and across
    resets. The package uses the decomposition in \cref{eq:additive-error} rather
    than collapsing every uncertainty into one white Gaussian sample.

    \begin{table}[htbp]
        \centering
        \caption{State and time-scale interpretation of additive error terms.}
        \label{tab:stochastic-terms}
        \begin{tabularx}{\textwidth}{@{}L{1.55in}L{1.5in}Y@{}}
\toprule
\textbf{Term} & \textbf{State lifecycle} & \textbf{Interpretation} \\
\midrule
Turn-on bias $\vect b_{\mathrm{on}}$ & Sampled once per reset & Run-to-run repeatability error; constant within a run. \\
White noise $\vect n_k$ & Independent each interval & Short-time rate scatter specified by density or continuous covariance. \\
Gauss--Markov bias $\vect b_k$ & Persistent, stationary & Correlated in-run drift with finite variance and correlation time. \\
Random-walk bias $\vect b_k$ & Persistent, nonstationary & Integrated driving noise whose variance grows with elapsed time. \\
Flicker-like bias $\vect f_k$ & Persistent, stationary sum & Finite-band approximation to a bias-instability plateau. \\
\bottomrule
        \end{tabularx}
    \end{table}

    This separation matters experimentally. A constant turn-on bias disappears from
    an Allan-deviation curve computed within one uninterrupted run, while a
    Gauss--Markov or flicker term remains visible through its time correlation.
    Conversely, a reset ensemble exposes turn-on repeatability directly.


    \section{Turn-on bias}
    For each channel,
    \begin{equation}
        \vect b_{\mathrm{on},s}\sim
        \Normal(\vect0,\sigma_{\mathrm{on},s}^{2}\mat I)
        \label{eq:turn-on-bias}
    \end{equation}
    is sampled independently at reset and held fixed. Conditional on one run it is
    a deterministic offset; across a population of resets it is a random variable.
    If the configured implementation supports axis-dependent or correlated
    repeatability, \(\sigma_{\mathrm{on}}^2\mat I\) can be replaced by a positive
    semidefinite covariance matrix without changing the lifecycle semantics.


    \section{White rate noise and sample-time scaling}
    Let \(\vect w_s(t)\) be continuous white rate noise with covariance intensity
    \(\mat Q_{c,s}\):
    \begin{equation}
        \E\!\left[\vect w_s(t)\vect w_s(t')^\transpose\right]
        =\mat Q_{c,s}\,\delta(t-t').
        \label{eq:white-continuous}
    \end{equation}
    The interval-average noise is
    \begin{equation}
        \vect n_{s,k}
        =\frac{1}{\Delta t_k}\int_{t_k}^{t_{k+1}}\vect w_s(t)\dd t,
    \end{equation}
    which has covariance
    \begin{align}
        \Cov(\vect n_{s,k})
        &=\frac{1}{\Delta t_k^2}
        \int_{t_k}^{t_{k+1}}\!\int_{t_k}^{t_{k+1}}
        \mat Q_{c,s}\delta(t-t')\dd t\dd t'\\
        &=\frac{\mat Q_{c,s}}{\Delta t_k}.
        \label{eq:white-discrete-cov}
    \end{align}
    Thus a scalar independent-axis density \(N_s\) produces
    \begin{equation}
        \vect n_{s,k}\sim
        \Normal\!\left(\vect0,\frac{N_s^2}{\Delta t_k}\mat I\right),
        \qquad
        \operatorname{std}(n_{s,k,i})=\frac{N_s}{\sqrt{\Delta t_k}}.
        \label{eq:white-density}
    \end{equation}
    After multiplying the rate by \(\Delta t_k\), the white-noise contribution to an
    increment has standard deviation
    \begin{equation}
        \operatorname{std}(\Delta y^{\mathrm{white}}_{s,k,i})
        =N_s\sqrt{\Delta t_k}.
        \label{eq:increment-noise-scaling}
    \end{equation}
    This rate-versus-increment distinction is a common source of factor-of-
    \(\sqrt{\Delta t}\) errors.


    \section{Thermal and cross-axis covariance}
    Define the thermal noise matrix
    \begin{equation}
        \mat\Lambda_{s,k}
        =\diag(\vect\lambda_{T,s,k}).
    \end{equation}
    The implementation uses
    \begin{equation}
        \vect n_{s,k}\sim
        \Normal\!\left(
                     \vect0,
                     \frac{\mat\Lambda_{s,k}\mat Q_{c,s}
                         \mat\Lambda_{s,k}^{\transpose}}{\Delta t_k}
        \right).
        \label{eq:general-white}
    \end{equation}
    When \field{noise\_covariance} is absent,
    \(\mat Q_{c,s}=N_s^2\mat I\). When it is supplied, that matrix replaces the
    independent-axis density path.

    A direct sampler can factor
    \begin{equation}
        \mat Q_{c,s}=\mat L_s\mat L_s^\transpose,
        \qquad
        \vect z_k\sim\Normal(\vect0,\mat I),
    \end{equation}
    and evaluate
    \begin{equation}
        \vect n_{s,k}
        =\frac{1}{\sqrt{\Delta t_k}}
        \mat\Lambda_{s,k}\mat L_s\vect z_k.
        \label{eq:correlated-sampling}
    \end{equation}
    A positive-semidefinite eigendecomposition is preferable to an ordinary
    Cholesky factor when a valid covariance is singular. Configuration validation
    should reject materially negative eigenvalues while allowing small
    floating-point residuals around zero.


    \section{Stationary first-order Gauss--Markov bias}
    A stationary first-order Gauss--Markov, or Ornstein--Uhlenbeck, process satisfies
    \begin{equation}
        \dd\vect b_s(t)
        =-\frac{1}{T_s}\vect b_s(t)\dd t
        +\sqrt{\frac{2\sigma_{b,s}^{2}}{T_s}}\dd\vect W(t),
        \label{eq:ou-sde}
    \end{equation}
    For axis-specific standard deviations, define \(\mat\Sigma_{b,s}=\diag(\vect\sigma_{b,s}\Had\vect\sigma_{b,s})\), with the scalar case
    represented by a multiple of the identity. The scalar equations below apply
    componentwise; their matrix form uses \(\mat\Sigma_{b,s}\).
    \begin{equation}
        \E\!\left[\vect b_s(t)\vect b_s(t+\tau)^\transpose\right]
        =\sigma_{b,s}^{2}e^{-\lvert\tau\rvert/T_s}\mat I.
        \label{eq:ou-autocov}
    \end{equation}

    Exact discretization over a potentially variable interval gives
    \begin{align}
        \vect b_{s,k+1}
        &=\phi_{s,k}\vect b_{s,k}+\vect\eta_{s,k},\\
        \phi_{s,k}&=e^{-\Delta t_k/T_s},\\
        \vect\eta_{s,k}
        &\sim\Normal\!\left(
                          \vect0,
            (1-\phi_{s,k}^{2})\mat\Sigma_{b,s}
        \right).
        \label{eq:gm-discrete}
    \end{align}
    Initializing
    \begin{equation}
        \vect b_{s,0}\sim\Normal(\vect0,\mat\Sigma_{b,s})
        \label{eq:gm-initial}
    \end{equation}
    starts the process in stationarity, avoiding an artificial warm-up transient.
    The discrete lag-\(\ell\) covariance for fixed \(\Delta t\) is
    \begin{equation}
        \Cov(\vect b_k,\vect b_{k+\ell})
        =\sigma_b^2\phi^{\lvert\ell\rvert}\mat I.
        \label{eq:gm-discrete-autocov}
    \end{equation}

    \begin{engineeringbox}
        The exact update in \cref{eq:gm-discrete} remains stable when the sample interval
        changes. An Euler update, by contrast, can distort both stationary variance and
        correlation when \(\Delta t/T\) is not very small.
    \end{engineeringbox}


    \section{Random-walk bias}
    When no correlation time is configured, the same scalar-or-three-axis bias
    parameter is interpreted as a random-walk driving standard deviation:
    \begin{align}
        \dd\vect b_s(t)&=\sigma_{\mathrm{rw},s}\dd\vect W(t),\\
        \vect b_{s,k+1}&=\vect b_{s,k}+\vect\eta_{s,k},\\
        \vect\eta_{s,k}&\sim
        \Normal\!\left(\vect0,
                     \sigma_{\mathrm{rw},s}^{2}\Delta t_k\mat I\right).
        \label{eq:random-walk}
    \end{align}
    The process is nonstationary:
    \begin{equation}
        \Var(b_i(t))=\Var(b_i(0))+\sigma_{\mathrm{rw},s}^{2}t.
        \label{eq:rw-variance}
    \end{equation}
    Therefore each component of \field{bias\_std} has different dimensional meaning
    depending on whether a correlation time is present. A profile and its
    documentation must state which interpretation applies.


    \section{Finite-band flicker-like bias}
    An ideal \(1/f\) process over an infinite frequency range is not a finite-variance
    stationary process. The package instead sums \(K\) independent stationary
    Gauss--Markov components,
    \begin{equation}
        \vect f_s(t)=\sum_{i=1}^{K}\vect x_{s,i}(t),
        \qquad
        T_i\in\operatorname{geomspace}(T_{\min},T_{\max},K),
        \label{eq:flicker-sum}
    \end{equation}
    with
    \begin{align}
        \vect x_{i,k+1}&=\phi_{i,k}\vect x_{i,k}+\vect\eta_{i,k},\\
        \phi_{i,k}&=e^{-\Delta t_k/T_i},\\
        \vect\eta_{i,k}&\sim\Normal\!\left(
                                         \vect0,s_i^2(1-\phi_{i,k}^2)\mat I\right).
        \label{eq:flicker-components}
    \end{align}
    The two-sided angular-frequency power spectral density of one component is
    \begin{equation}
        S_{x_i}(\omega)=\frac{2s_i^2T_i}{1+\omega^2T_i^2}.
        \label{eq:ou-psd}
    \end{equation}
    A logarithmic bank of such Lorentzian spectra can approximate a
    \(1/\lvert\omega\rvert\) region between characteristic frequencies on the order
    of \(1/T_{\max}\) and \(1/T_{\min}\). Outside that finite band, the spectrum
    rolls off rather than remaining ideal flicker noise.

    Let \(\mathcal A(v,T,\tau)\) denote the Allan variance of a stationary
    Gauss--Markov component with stationary variance \(v\):
    \begin{equation}
        \mathcal A(v,T,\tau)=
        \frac{
            2v\left[T\tau-T^2(1-e^{-\tau/T})\right]
            -vT^2(1-e^{-\tau/T})^2
        }{\tau^2}.
        \label{eq:gm-allan-function}
    \end{equation}
    At the geometric-center averaging time
    \begin{equation}
        \tau_c=\sqrt{T_{\min}T_{\max}},
    \end{equation}
    the implementation uses equal component standard deviations
    \begin{equation}
        s_i=\frac{\sigma_f}{
            \sqrt{\displaystyle\sum_{\ell=1}^{K}
                \mathcal A(1,T_\ell,\tau_c)}},
        \qquad
        \sigma_f=\field{flicker\_bias\_std}.
        \label{eq:flicker-normalization}
    \end{equation}
    Consequently, \field{flicker\_bias\_std} is the target Allan deviation of the
    \emph{sum} at \(\tau_c\), not the stationary standard deviation of one component.
    Increasing \field{flicker\_components} smooths the approximation but does not
    widen its configured time-constant band.


    \section{Combined covariance and conditioning on a run}
    Before clipping and quantization, and conditional on the persistent states, the
    only new random term at interval \(k\) is white noise. Across a reset ensemble,
    the instantaneous additive error covariance can be decomposed as
    \begin{equation}
        \Cov(\vect e_{s,k})
        =\mat P_{\mathrm{on},s}
        +\mat P_{b,s,k}
        +\mat P_{f,s,k}
        +\frac{\mat\Lambda_{s,k}\mat Q_{c,s}
            \mat\Lambda_{s,k}^\transpose}{\Delta t_k},
        \label{eq:combined-cov}
    \end{equation}
    provided the terms are independent. Within one initialized run,
    \(\vect b_{\mathrm{on},s}\) is fixed and contributes to the conditional mean,
    not to interval-to-interval conditional covariance. This distinction should be
    preserved in uncertainty budgets and tests.


    \section{Random-number ownership and reproducibility}
    A stateful model should own or receive an explicit pseudo-random generator. A
    fixed seed should reproduce the full reset state and subsequent sample sequence
    for an identical call history. Reproducibility requires more than reseeding
    white noise: turn-on bias, misalignment, Gauss--Markov initialization, and every
    flicker component must draw from the same defined lifecycle.

    Changing the sample schedule changes both the number of random draws and the
    exact transition covariance. Therefore two runs with the same seed but different
    sampling times are not expected to share pointwise values; they should share the
    same configured continuous-time statistics.


    \chapter{Sampling algorithm and state lifecycle}
    \label{chap:algorithm}


    \section{State carried by one model instance}
    A model instance is logically associated with one sensor stream. Its mutable
    state includes
    \begin{itemize}
        \item the previous timestamp, velocity, and orientation;
        \item reset-level turn-on bias and misalignment for each channel;
        \item the current Gauss--Markov or random-walk state;
        \item all finite-band flicker component states; and
        \item pseudo-random generator state.
    \end{itemize}
    Sharing one instance between unrelated vehicles or asynchronous streams would
    mix their histories. Use one instance per logical IMU unless a specialized
    implementation explicitly supports multiplexing.


    \section{Reference interval algorithm}
    \Cref{alg:measure} states the implemented lifecycle. The persistent bias and
    flicker states are advanced using the current interval duration before they are
    added to that interval's returned rate; the updated states therefore correspond
    to the interval ending at the current timestamp.

    \begin{algorithm}[htbp]
\caption{Reference \software{measure} operation for one interval.}
\label{alg:measure}
\DontPrintSemicolon
\KwIn{$t$, $\vect v^w$, $\Rwb$, optional $T_{\mathrm C}$}
\KwOut{Measured $\Delta\vect v$, $\Delta\vect\theta$, interval metadata}
Validate finite inputs and proper rotation matrix\;
\eIf{history is uninitialized}{
    Store $t$, $\vect v^w$, and $\Rwb$\;
    Return zero increments with a zero-length initialization result\;
}{
    Set $\Delta t\leftarrow t-t_{\mathrm{prev}}$ and require $\Delta t>0$\;
    Compute ideal increments using \cref{eq:ideal-dv,eq:ideal-dtheta}\;
    Compute ideal channel rates using \cref{eq:ideal-rates}\;
    \For{$s\in\{a,g\}$}{
        Evaluate thermal state at $T_{\mathrm C}$\;
        Form deterministic rate $\vect d_s$ and rotate by fixed $\mat M_s$\;
        Advance Gauss--Markov/random-walk and flicker states using $\Delta t$\;
        Sample white noise for this $\Delta t$\;
        Add turn-on, in-run bias, flicker, and white terms\;
        Apply clipping, quantization, output scaling, and integration\;
    }
    Store current truth as previous truth\;
    Return both measured increments and $[t_{\mathrm{prev}},t]$ metadata\;
}
    \end{algorithm}

    This ordering is part of the runtime contract: changing it changes seeded
    outputs even when the stationary distributions remain the same. Equations,
    tests, and saved seeds must use the same state-timing convention.


    \section{Numerical considerations}

    \subsection{Small rotations}
    For \(\lVert\vect\phi\rVert\ll1\), use
    \begin{align}
        \frac{\sin\theta}{\theta}
        &\approx1-\frac{\theta^2}{6}+\frac{\theta^4}{120},\\
        \frac{1-\cos\theta}{\theta^2}
        &\approx\frac12-\frac{\theta^2}{24}+\frac{\theta^4}{720}
    \end{align}
    in \cref{eq:rodrigues}. This avoids cancellation and preserves smooth behavior
    at zero rotation.

    \subsection{Correlation-time extremes}
    When \(\Delta t/T\) is very small, evaluate
    \begin{equation}
        1-e^{-2\Delta t/T}
        =-\operatorname{expm1}(-2\Delta t/T)
    \end{equation}
    with an \software{expm1}-style function to avoid subtracting nearly equal
    numbers. When \(\Delta t/T\) is very large, \(\phi\) tends to zero and the next
    state becomes an independent draw from the stationary distribution.

    \subsection{Quantization and floating-point ties}
    The exact tie-breaking behavior of \(\round\) should be documented because
    banker's rounding and away-from-zero rounding differ at half-step inputs.
    Regression fixtures should include positive and negative half-step cases.

    \subsection{Temperature absence}
    \software{None} is a semantic value meaning that temperature is unavailable; it
    is not equivalent to \SI{0}{\degreeCelsius}. The output should preserve that
    absence rather than inventing a temperature sample.


    \section{Reset semantics}
    A reset clears truth history and samples fresh run-level stochastic quantities.
    The first post-reset measurement again initializes history and returns zero
    increments. A deterministic test can compare two resets with the same seed;
    a variability test can compare many seeds and verify the configured reset
    ensemble covariance.


    \chapter{Allan deviation and spectral interpretation}
    \label{chap:allan}


    \section{Estimator used by the analysis examples}
    For equally spaced rate samples \(y_i\) with sample period \(\Delta t\), the
    analysis uses non-overlapping blocks of \(m\) samples. With \(J\) complete block
    means,
    \begin{align}
        \tau&=m\Delta t,\\
        \bar y_j^{(m)}
        &=\frac{1}{m}\sum_{i=jm}^{(j+1)m-1}y_i,\\
        \widehat\sigma_A^2(\tau)
        &=\frac{1}{2(J-1)}
        \sum_{j=0}^{J-2}
        \left(\bar y_{j+1}^{(m)}-\bar y_j^{(m)}\right)^2.
        \label{eq:allan-estimator}
    \end{align}
    Only averaging times \(\tau=m\Delta t\) with at least two adjacent complete
    blocks are defined. As \(\tau\) grows, \(J\) decreases and statistical
    uncertainty rises. Allan analysis and its interpretation are described in
    ~\cite{allan1966,riley2008,elsheimy2008}.


    \section{Reference curves}
    For independent white rate noise with density \(N\),
    \begin{equation}
        \sigma_A(\tau)=\frac{N}{\sqrt\tau},
        \label{eq:allan-white}
    \end{equation}
    which has slope \(-1/2\) on log--log axes. For a rate random walk driven by
    \(\sigma_{\mathrm{rw}}\),
    \begin{equation}
        \sigma_A(\tau)=\sigma_{\mathrm{rw}}\sqrt{\frac{\tau}{3}},
        \label{eq:allan-rw}
    \end{equation}
    with slope \(+1/2\).

    For the stationary Gauss--Markov process in \cref{eq:ou-sde},
    \begin{equation}
        \sigma_A^2(\tau)=\sigma_b^2\left[
                                       \frac{2T}{\tau}
                                       -\left(\frac{T}{\tau}\right)^2
                                       \left(3-4e^{-\tau/T}+e^{-2\tau/T}\right)
        \right].
        \label{eq:allan-gm}
    \end{equation}
    Its limiting behavior is
    \begin{align}
        \sigma_A(\tau)
        &\sim\sigma_b\sqrt{\frac{2\tau}{3T}},
        &&\tau\ll T,\label{eq:gm-allan-small}\\
        \sigma_A(\tau)
        &\sim\sigma_b\sqrt{\frac{2T}{\tau}},
        &&\tau\gg T.
        \label{eq:gm-allan-large}
    \end{align}
    The process therefore rises with slope \(+1/2\), turns near its correlation
    time, and falls with slope \(-1/2\).

    \begin{figure}[htbp]
        \centering
        \begin{tikzpicture}
            \begin{loglogaxis}
                [
                width=0.88\textwidth,
                height=0.49\textwidth,
                xmin=0.02,xmax=100,
                ymin=0.07,ymax=8,
                xlabel={Normalized averaging time $\tau/T$},
                ylabel={Normalized Allan deviation},
                grid=both,
                major grid style={draw=ImuRule},
                minor grid style={draw=ImuRule!55},
                tick label style={font=\small},
                label style={font=\small},
                legend style={at={(0.5,-0.22)},anchor=north,legend columns=3,draw=none,font=\small}
                ]
                \addplot[ImuBlue,very thick,domain=0.02:100,samples=180] {x^(-0.5)};
                \addlegendentry{white rate noise}
                \addplot[ImuOrange,very thick,dashed,domain=0.02:100,samples=180] {sqrt(x/3)};
                \addlegendentry{rate random walk}
                \addplot[ImuTeal,very thick,domain=0.02:100,samples=220]
                {sqrt(2/x - (1/x)^2*(3 - 4*exp(-x) + exp(-2*x)))};
                \addlegendentry{stationary Gauss--Markov}
            \end{loglogaxis}
        \end{tikzpicture}
        \caption{Normalized Allan-deviation signatures. Absolute levels depend on
            $N$, $\sigma_{\mathrm{rw}}$, and $\sigma_b$; the plotted shapes emphasize slope
            and turnover rather than a particular sensor.}
        \label{fig:allan-reference}
    \end{figure}

    \begin{table}[htbp]
        \centering
        \caption{Expected Allan-deviation signatures for configured model terms.}
        \label{tab:allan-signatures}
        \begin{tabularx}{\textwidth}{@{}L{1.55in}L{1.65in}Y@{}}
\toprule
\textbf{Term} & \textbf{Signature} & \textbf{Important interpretation} \\
\midrule
White rate noise & $N/\sqrt\tau$; slope $-1/2$ & Verifies density and sample-time scaling. \\
Rate random walk & $\sigma_{\mathrm{rw}}\sqrt{\tau/3}$; slope $+1/2$ & Nonstationary long-time growth. \\
Gauss--Markov & Rise, turnover near $T$, then fall & Both stationary standard deviation and correlation time matter. \\
Finite-band flicker sum & Approximately flat over configured band & \field{flicker\_bias\_std} targets the sum at $\tau_c$. \\
Turn-on bias & No within-run contribution & Must be assessed across resets, not one Allan trace. \\
Clipping / quantization & Signal- and level-dependent & Can break simple Gaussian power-law interpretations. \\
\bottomrule
        \end{tabularx}
    \end{table}


    \section{Addition of independent Allan variances}
    If independent zero-mean processes sum,
    \begin{equation}
        y(t)=\sum_j y_j(t),
    \end{equation}
    then their Allan variances add:
    \begin{equation}
        \sigma_{A,y}^2(\tau)=\sum_j\sigma_{A,y_j}^2(\tau).
        \label{eq:allan-addition}
    \end{equation}
    Their Allan \emph{deviations} do not add directly. This distinction is used in
    \cref{eq:flicker-normalization} and should also be used when plotting analytic
    reference envelopes for a composite configuration.


    \section{Finite records and confidence}
    For an idealized variance estimate with effective degrees of freedom \(\nu\), an
    approximate two-sided \((1-\alpha)\) confidence interval is
    \begin{equation}
        \frac{\nu\widehat\sigma_A^2}
        {\chi^2_{1-\alpha/2,\nu}}
        \le \sigma_A^2\le
        \frac{\nu\widehat\sigma_A^2}
        {\chi^2_{\alpha/2,\nu}}.
        \label{eq:allan-confidence}
    \end{equation}
    The effective \(\nu\) depends on estimator overlap and noise type, so
    \cref{eq:allan-confidence} is a planning relation rather than a universal drop-in
    formula. Regression tests should use tolerances informed by record length and
    Monte Carlo scatter instead of requiring every empirical point to lie exactly on
    an analytic curve.


    \section{Bias-instability naming}
    Datasheets and Allan-analysis references sometimes report a bias-instability
    coefficient \(B\) related to a flat Allan-deviation level by a convention such as
    \begin{equation}
        \sigma_A\approx0.664 B.
    \end{equation}
    The package avoids silently imposing that convention:
    \field{flicker\_bias\_std} is defined directly as an Allan-deviation target at
    \(\tau_c\). A profile importer should record whether a source value is already an
    Allan level or is a conventionally normalized coefficient before conversion.


    \chapter{Downstream error growth and engineering interpretation}
    \label{chap:propagation}


    \section{Why separate error terms matter}
    The package does not propagate a navigation solution, but the downstream effects
    of its error terms explain why each parameter deserves an independent model.
    For a conventional strapdown mechanization, a simplified linearized error system
    has the form
    \begin{align}
        \delta\dot{\vect\theta}
        &\approx-\vect b_g-\vect n_g,
        \label{eq:attitude-error-dynamics}\\
        \delta\dot{\vect v}
        &\approx-\mat R\skewmat{\vect f}\,\delta\vect\theta
        +\mat R(\vect b_a+\vect n_a),
        \label{eq:velocity-error-dynamics}\\
        \delta\dot{\vect p}&=\delta\vect v,
        \label{eq:position-error-dynamics}
    \end{align}
    where signs and detailed transport terms depend on the chosen navigation-frame
    convention. These relations are analytical context, not computations performed
    by \pkg{}~\cite{titterton2004,groves2013,woodman2007}.


    \section{Constant-bias approximations}
    Ignoring frame rotation and feedback, a constant accelerometer bias produces
    \begin{align}
        \delta\vect v(t)&\approx\vect b_a t,\\
        \delta\vect p(t)&\approx\frac12\vect b_a t^2.
        \label{eq:accel-bias-growth}
    \end{align}
    A constant gyroscope bias produces attitude error magnitude approximately
    \begin{equation}
        \lVert\delta\vect\theta(t)\rVert
        \approx\lVert\vect b_g\rVert t.
        \label{eq:gyro-bias-attitude}
    \end{equation}
    Near level flight, that tilt misprojects gravity into horizontal acceleration:
    \begin{equation}
        \delta\vect a_h(t)\approx\vect g\times\delta\vect\theta(t).
        \label{eq:tilt-gravity}
    \end{equation}
    For a constant gyro bias and fixed gravity direction, the corresponding leading
    position term scales cubically,
    \begin{equation}
        \delta\vect p_h(t)
        \approx\frac16\left(\vect g\times\vect b_g\right)t^3,
        \label{eq:gyro-bias-position}
    \end{equation}
    up to sign and frame convention. This mechanism is one reason small gyro errors
    can dominate unaided inertial drift.


    \section{White-noise covariance growth}
    For one axis driven by continuous white accelerometer noise density \(N_a\),
    \begin{align}
        \Var(\delta v(t))&=N_a^2t,\\
        \Var(\delta p(t))&=\frac{N_a^2t^3}{3}.
        \label{eq:accel-white-growth}
    \end{align}
    For continuous white gyro rate noise density \(N_g\), direct attitude-error
    variance grows as
    \begin{equation}
        \Var(\delta\theta(t))=N_g^2t.
        \label{eq:gyro-white-growth}
    \end{equation}
    These formulas assume ideal continuous integration and no estimator correction,
    but they provide useful order-of-magnitude checks for Monte Carlo experiments.


    \section{Run-to-run versus time-varying uncertainty}
    Suppose a simulation campaign runs the same trajectory repeatedly. A turn-on
    bias shifts each run coherently, while white noise varies sample by sample. The
    law of total covariance separates these effects:
    \begin{equation}
        \Cov(\vect y)
        =\E\!\left[\Cov(\vect y\mid\vect b_{\mathrm{on}})\right]
        +\Cov\!\left(\E[\vect y\mid\vect b_{\mathrm{on}}]\right).
        \label{eq:total-covariance}
    \end{equation}
    The first term describes within-run scatter and the second run-to-run
    repeatability. A validation campaign that never resets cannot estimate the
    second term.

    \begin{table}[htbp]
        \centering
        \caption{Leading unaided error-growth intuition under simplified assumptions.}
        \label{tab:error-growth}
        \begin{tabularx}{\textwidth}{@{}L{1.65in}L{1.55in}Y@{}}
\toprule
\textbf{Sensor error} & \textbf{Primary state effect} & \textbf{Leading growth} \\
\midrule
Constant accelerometer bias & Velocity, then position & $\delta v\propto t$, $\delta p\propto t^2$ \\
Accelerometer white noise & Velocity random walk & $\operatorname{std}(\delta v)\propto t^{1/2}$; $\operatorname{std}(\delta p)\propto t^{3/2}$ \\
Constant gyro bias & Attitude, gravity projection & $\delta\theta\propto t$; horizontal $\delta p\propto t^3$ in a level simplification \\
Gyro white noise & Attitude random walk & $\operatorname{std}(\delta\theta)\propto t^{1/2}$ \\
Scale factor & Signal-proportional error & Grows with maneuver magnitude and duration \\
Misalignment & Cross-axis coupling & Depends on signal orthogonal to the affected axis \\
Saturation & Information loss & Potentially unbounded downstream effect during clipped maneuvers \\
\bottomrule
        \end{tabularx}
    \end{table}


    \section{Implications for controller and estimator tests}
    A useful campaign should not stop at a zero-motion Allan run. It should include
    \begin{itemize}
        \item static tests that isolate bias and white noise;
        \item single-axis ramps that expose scale factor and quantization;
        \item multi-axis rotations that expose misalignment and nonlinear coupling;
        \item temperature sweeps that separate mean drift from noise-level change;
        \item high-dynamic cases that exercise saturation; and
        \item closed-loop scenarios that compare truth-fed, ideal-IMU, and corrupted-IMU
        performance.
    \end{itemize}
    This layered approach identifies whether a failure originates in vehicle
    dynamics, guidance and control, state estimation, or sensor corruption.


    \chapter{Configuration, units, and profile construction}
    \label{chap:configuration}


    \section{Canonical-unit policy}
    The runtime model expects configuration values in the canonical units below. The
    loader validates values but does not perform unit conversion; conversion belongs
    at the profile boundary, not deep inside stochastic updates. Temperature is an
    explicit Celsius boundary quantity because the implemented thermal model uses
    temperature differences. The recommended canonical units are shown in
    \cref{tab:config-units}.

    \begin{table}[htbp]
        \centering
        \caption{Recommended canonical units for model parameters.}
        \label{tab:config-units}
        \begin{tabularx}{\textwidth}{@{}L{1.7in}L{2.0in}Y@{}}
\toprule
\textbf{Parameter family} & \textbf{Accelerometer} & \textbf{Gyroscope} \\
\midrule
Rate and bias & \si{\meter\per\second\squared} & \si{\radian\per\second} \\
White-noise density & \(\si{\meter\per\second\squared}/\sqrt{\si{\hertz}}\) & \(\si{\radian\per\second}/\sqrt{\si{\hertz}}\) \\
Continuous covariance & $(\si{\meter\per\second\squared})^2\!\cdot\si{\second}$ & $(\si{\radian\per\second})^2\!\cdot\si{\second}$ \\
Random-walk drive & \(\si{\meter\per\second\squared}/\sqrt{\si{\second}}\) & \(\si{\radian\per\second}/\sqrt{\si{\second}}\) \\
Scale factor / nonlinearity gain & dimensionless / inverse rate squared & dimensionless / inverse rate squared; scale factor may be three-axis \\
Misalignment & \si{\radian} or radians squared covariance & \si{\radian} or radians squared covariance \\
Range and quantization & \si{\meter\per\second\squared} & \si{\radian\per\second}; quantization may be three-axis \\
Thermal bias coefficient & \si{\meter\per\second\squared\per\degreeCelsius} & \si{\radian\per\second\per\degreeCelsius} \\
\bottomrule
        \end{tabularx}
    \end{table}


    \section{Common conversions}
    The following conversions are useful when importing datasheet or laboratory
    values:
    \begin{align}
        1\ \si{\degree\per\hour}
        &=\frac{\pi}{180\cdot3600}\ \si{\radian\per\second},\\
        1\ \si{\degree}/\sqrt{\si{\hour}}
        &=\frac{\pi}{180\cdot60}\ \si{\radian}/\sqrt{\si{\second}},\\
        1\ \si{\micro\gzero}/\sqrt{\si{\hertz}}
        &=10^{-6}g_0\ \si{\meter\per\second\squared}/\sqrt{\si{\hertz}},\\
        1\ \si{\ppm}&=10^{-6},\\
        1\ \si{\arcsec}&=\frac{\pi}{180\cdot3600}\ \si{\radian}.
        \label{eq:unit-conversions}
    \end{align}
    Use one documented value of standard gravity, conventionally
    \(g_0=\SI{9.80665}{\meter\per\second\squared}\), throughout a profile conversion.


    \section{Mapping source specifications to model terms}
    \begin{table}[htbp]
        \centering
        \caption{Typical source metric to configuration mapping.}
        \label{tab:profile-mapping}
        \begin{tabularx}{\textwidth}{@{}L{1.8in}L{1.8in}Y@{}}
\toprule
\textbf{Source metric} & \textbf{Model term} & \textbf{Caution} \\
\midrule
Noise density / angle or velocity random walk & $N_s$ or $\mat Q_{c,s}$ & Confirm whether the source uses rate samples, integrated angle/velocity, one-sided PSD, or Allan normalization. \\
Bias repeatability / turn-on bias & $\sigma_{\mathrm{on},s}$ & Usually a reset ensemble metric, not an in-run drift metric. \\
Bias stability / in-run stability & GM and/or flicker parameters & One scalar rarely determines both stationary variance and correlation time; fit a curve when possible. \\
Scale-factor error & $\vect s_s$ & Convert percent or ppm to dimensionless gain. Distinguish initial error from thermal coefficient and preserve axis dependence when available. \\
Nonlinearity & $n_s$ & The package uses a radial quadratic gain; a vendor's definition may use full-scale percentage or a different polynomial. \\
Axis alignment & $\mat\Sigma_{\epsilon,s}$ or a fixed calibration matrix & Convert angular units to radians and determine whether the source reports a rotation-vector covariance, one-sigma isotropic value, maximum, or calibrated residual. \\
Measurement range & $L_s$ & Determine whether the published range is symmetric and whether clipping occurs before or after internal filtering. \\
Digital resolution / LSB & $\vect\delta_s$ & Resolution, noise-free counts, and effective number of bits are not interchangeable; preserve per-axis LSBs when available. \\
Temperature coefficient & $\vect c_b$, $\vect c_n$, $\vect c_{\mathrm{sf}}$ & Verify reference temperature, axis dependence, and whether the relation is linear over the intended range. \\
\bottomrule
        \end{tabularx}
    \end{table}

    \begin{warningbox}
        Do not infer a Gauss--Markov correlation time from a bias-stability number alone.
        Many pairs \((\sigma_b,T)\) can produce similar Allan levels over a limited
        record. Preserve the source metric, fitting method, and assumptions in profile
        metadata.
    \end{warningbox}


    \section{Profile provenance}
    The format-neutral \software{load\_profile\_document} loader validates both
    configuration and the descriptive metadata schema. It records profile
    classification, review dates, source URLs, and optional archive links. For
    unambiguous interpretation and reproducibility, profile provenance should also
    preserve
    \begin{itemize}
        \item profile identifier and schema version;
        \item source citations, page/table references, and retrieval date;
        \item whether each value is authoritative, fitted, inferred, or purely
        notional;
        \item original value and units before conversion;
        \item canonical value and conversion rule;
        \item uncertainty convention: one-sigma, root-mean-square, maximum, typical,
        or unknown;
        \item applicable temperature, bandwidth, sample rate, and operating mode;
        \item assumptions used to map source metrics into model parameters; and
        \item license or redistribution constraints for embedded source material.
    \end{itemize}

    A useful evidence tier is
    \begin{enumerate}
        \item \textbf{Traceable source profile}: direct values with authoritative
        citation and minimal interpretation.
        \item \textbf{Fitted engineering profile}: parameters fitted to supplied
        measurements or Allan data with documented residuals.
        \item \textbf{Proxy profile}: engineering estimate based on a comparable
        device or class.
        \item \textbf{Test fixture}: intentionally synthetic values chosen to isolate
        one implementation behavior.
    \end{enumerate}
    The repository's hardware-oriented examples are notional unless their metadata
    establishes otherwise.


    \section{Positive-semidefinite covariance construction}
    For per-axis densities \(N_x,N_y,N_z\) and a correlation matrix
    \(\mat P\), define
    \begin{equation}
        \mat D_N=\diag(N_x,N_y,N_z),
        \qquad
        \mat Q_c=\mat D_N\mat P\mat D_N.
        \label{eq:covariance-construction}
    \end{equation}
    A valid correlation matrix is symmetric, has unit diagonal, and is positive
    semidefinite. Pairwise correlations inside \([-1,1]\) are necessary but not
    sufficient for a valid three-axis matrix; validate eigenvalues after assembly.


    \section{Worked conversion example}
    Consider an illustrative gyro source reporting
    \begin{equation}
        N_g=0.20\ \si{\degree}/\sqrt{\si{\hour}},
        \qquad
        \sigma_{\mathrm{on}}=5\ \si{\degree\per\hour},
        \qquad
        s_g=500\ \si{\ppm}.
    \end{equation}
    The canonical values are
    \begin{align}
        N_g
        &=0.20\frac{\pi}{180\cdot60}
        \approx5.82\times10^{-5}\,\si{\radian}/\sqrt{\si{\second}},\\
        \sigma_{\mathrm{on}}
        &=5\frac{\pi}{180\cdot3600}
        \approx\SI{2.42e-5}{\radian\per\second},\\
        s_g&=500\times10^{-6}=5.0\times10^{-4}.
    \end{align}
    These numbers are pedagogical only. A real profile must retain the source's
    statistical conventions and operating conditions.


    \chapter{Software contract and integration guidance}
    \label{chap:software}


    \section{Public concepts}
    The source contract names four central concepts:
    \begin{description}
        \item[\software{ImuModelProtocol}]
        A structural interface for any stateful implementation that accepts truth
        samples and returns IMU increments.
        \item[\software{CheckpointableImuModelProtocol}]
        An optional generic structural interface exposing model-defined snapshot and
        restore operations for implementations that can pause and resume their
        stochastic state exactly. It does not prescribe JSON, Pydantic, or disk-file serialization.
        \item[\software{SerializableCheckpointableImuModelProtocol}]
        An optional refinement that exposes a typed byte codec for a checkpointable
        model. Its codec may use JSON, MessagePack, CBOR, or another transport format.
        \item[\software{ImuConfig}]
        Top-level configuration, including global output scaling and channel
        configuration.
        \item[\software{ImuOutput}]
        Output carrying \field{delta\_v}, \field{delta\_theta}, interval timestamps,
        derived \field{dt}, and \field{temperature\_celsius} when available.
        \item[Thermal model]
        A configured channel object exposing
        \software{evaluate(temperature\_celsius)}. Coefficients are bound at
        construction, not passed on every evaluation.
        \item[\software{ImuErrorProcessEstimate}]
        A stateless, estimator-facing characterization of configured sensor noise,
        persistent-error transitions, priors, and increment covariances.
    \end{description}
    Alternate implementations satisfy \software{ImuModelProtocol} through structural
    typing, allowing high-fidelity, hardware-specific, or accelerated models to
    replace the default without changing consumers.


    \section{Conceptual integration}
    The following listing is intentionally schematic; exact import paths and
    constructors belong in generated API documentation. It emphasizes the lifecycle
    and frame contract.

    \begin{lstlisting}[style=imuPython,caption={Conceptual integration with a truth-producing simulation.},label={lst:integration}]
model = configured_imu_model
model.reset()  # samples run-level bias and misalignment

# The first call establishes history and returns zero increments.
output_0 = model.measure(
    timestamp=t_0,
    velocity_without_gravity=velocity_world_0,
    orientation_world_from_body=rotation_world_from_body_0,
    temperature_celsius=temperature_0,
)

# Each later call returns the increment over the preceding interval.
output_1 = model.measure(
    timestamp=t_1,
    velocity_without_gravity=velocity_world_1,
    orientation_world_from_body=rotation_world_from_body_1,
    temperature_celsius=temperature_1,
)

estimator.propagate(
    delta_v_body_start=output_1.delta_v,
    delta_theta_body_start=output_1.delta_theta,
    dt=output_1.dt,
)
    \end{lstlisting}


    \section{Placement in a time-stepped simulation}
    For a controller that must act on sensed rather than truth state, one useful step
    ordering is
    \begin{enumerate}
        \item the equations of motion expose the truth snapshot at the sampling epoch;
        \item the IMU error model converts the previous and current snapshots into
        corrupted increments;
        \item the estimator updates its internal state from those increments and any
        other measurements;
        \item guidance and control operate on the estimated state; and
        \item the equations of motion advance using the commanded actuation.
    \end{enumerate}
    This ordering prevents a controller from inadvertently seeing end-of-step truth
    before it makes the decision associated with that step. The exact scheduler may
    use begin-step or end-step sampling, but the timestamp and interval convention
    must remain consistent.

    \begin{figure}[htbp]
        \centering
        \begin{tikzpicture}[
            node distance=8mm and 13mm,
            block/.style={draw=ImuRule,rounded corners=2pt,fill=white,
            minimum height=11mm,text width=33mm,align=center,font=\small},
            arrow/.style={-{Latex[length=2.5mm]},line width=0.85pt,draw=ImuBlue}
        ]
  \node[block,fill=ImuPale] (eom) {Equations of motion\\truth at $t_k$};
  \node[block,right=of eom,fill=ImuPaleTeal] (imu) {IMU error model\\$\Delta\vect v,\Delta\vect\theta$};
  \node[block,right=of imu,fill=ImuPale] (est) {Estimator\\sensed state};
  \node[block,below=of est,fill=ImuPaleGold] (ctl) {Guidance / controller\\command};
  \node[block,left=of ctl,fill=ImuPale] (act) {Actuator / plant input};
  \draw[arrow] (eom) -- (imu);
  \draw[arrow] (imu) -- (est);
  \draw[arrow] (est) -- (ctl);
  \draw[arrow] (ctl) -- (act);
  \draw[arrow] (act) -| (eom);
        \end{tikzpicture}
        \caption{Recommended closed-loop placement: truth is converted to sensor
        increments before estimator and controller logic execute.}
        \label{fig:closed-loop}
    \end{figure}


    \section{State ownership and concurrency}
    Because measurement history and stochastic processes are mutable, a model
    instance is not a pure function of the current sample. Consumers should
    \begin{itemize}
        \item create one instance per logical sensor stream;
        \item serialize calls in timestamp order;
        \item avoid sharing an instance between Monte Carlo trials;
        \item reset explicitly between runs; and
        \item capture configuration, seed, and sample schedule in reproducibility
        metadata.
    \end{itemize}
    Thread safety should not be assumed unless a concrete implementation documents
    its synchronization behavior.


    \section{Checkpoint and resume}
    The built-in \software{ImuModel} supports an explicit pause/resume boundary.
    \software{snapshot()} returns the versioned Pydantic model
    \software{ImuModelCheckpoint}; \software{save\_checkpoint(path)} uses the
    built-in JSON codec to write that envelope as UTF-8, and
    \software{load\_checkpoint(path)} or \software{from\_checkpoint(checkpoint)}
    reconstructs a model ready for its next measurement.

    A checkpoint taken between completed \software{measure} calls contains the
    canonical configuration, NumPy bit-generator class and state, turn-on bias,
    current Gauss--Markov/random-walk bias, every finite-band flicker component,
    sampled misalignment, and the previous timestamp, velocity, and orientation.
    Restoring all of these values is necessary: restoring only a seed would not
    reproduce a run after the random stream had advanced, and restoring only the
    latent biases would lose the interval baseline.

    \begin{lstlisting}[style=imuPython,caption={Pause and resume a stochastic model.},label={lst:checkpoint}]
model.save_checkpoint("run-checkpoint.json")

# The file can be moved to another process or host with the same package contract.
resumed = ImuModel.load_checkpoint("run-checkpoint.json")
output = resumed.measure(
    timestamp=t_next,
    velocity_without_gravity=velocity_world_next,
    orientation_world_from_body=rotation_world_from_body_next,
    temperature_celsius=temperature_next,
)
    \end{lstlisting}

    The built-in codec is format-specific and uses JSON rather than a Python
    pickle. Its byte boundary is suitable for files or wire transport and its
    payload includes a schema version and model identifier so unsupported data
    fails explicitly. The generic \software{CheckpointableImuModelProtocol}
    requires only model-defined snapshot and restore operations; the optional
    \software{SerializableCheckpointableImuModelProtocol} adds a typed codec
    for byte transport. A different implementation may use a different checkpoint
    type and a MessagePack, CBOR, or binary codec. A custom \software{ThermalModel}
    may carry additional state and therefore must provide its own persistence
    mechanism before it can participate in an exact checkpoint. File writes are
    atomic, and \software{save\_checkpoint(..., codec=...)} and
    \software{load\_checkpoint(..., codec=...)} allow an explicitly selected
    codec when JSON is not the desired transport representation.


    \section{Extension points}
    A custom model can preserve the public contract while adding, for example,
    \begin{itemize}
        \item a measured frequency-response or anti-aliasing filter;
        \item full calibration matrices rather than scalar scale factors;
        \item cross-channel accelerometer--gyro correlation;
        \item temperature-state dynamics and thermal lag;
        \item vibration rectification, gyro g-sensitivity, or shock recovery;
        \item coning and sculling corrections using substep truth; or
        \item hardware replay and hybrid measured/simulated channels.
    \end{itemize}
    Such features should be isolated behind the protocol rather than weakening the
    meaning of existing fields.


    \chapter{Stateless error-process characterization}
    \label{chap:characterization}

    The stateful model is appropriate when a simulation needs a realized stochastic
    trajectory. Estimators and covariance studies need a different contract: they
    need to know what the configured sensor process says is statistically plausible,
    without advancing a model instance or drawing a particular bias realization.
    The public function
    \software{characterize\_imu\_error\_process} provides that contract.

    \begin{tcolorbox}[
        enhanced,breakable,colback=ImuPaleTeal,colframe=ImuTeal,
        title={Pure characterization contract},fonttitle=\bfseries]
        \begin{lstlisting}[style=imuPython]
if imu_output.dt > 0.0:
    noise = characterize_imu_error_process(
        config=imu_config,
        dt=imu_output.dt,
        process_age_at_interval_start_s=(
            imu_output.start_time - initialization_time
        ),
        temperature_celsius=imu_output.temperature_celsius,
        include_quantization_approximation=False,
    )
        \end{lstlisting}
        The function performs no random-number generation, does not inspect a live
        \software{ImuModel}, does not depend on a realized turn-on or current bias, and
        does not mutate its configuration. Equal inputs produce equal statistics.
        The function requires finite \software{dt > 0} and
        \software{process\_age\_at\_interval\_start\_s >= 0}. The zero-duration output
        from the first \software{measure} call is an initialization sentinel and is not a
        valid characterization interval. Temperature is explicitly in degrees Celsius.
    \end{tcolorbox}


    \section{Process-age and state-timing convention}
    Let $a_k\ge0$ be the modeled process age at the start of an interval and let
    $a_{k+1}=a_k+\Delta t_k$. The age origin is the first truth sample after
    construction or reset, because that call initializes history without advancing
    in-run stochastic states; it is not necessarily wall-clock time since
    \software{reset()} was called.

    The runtime sampler advances persistent bias and flicker states before applying
    them to the returned rate. Consequently,
    \begin{equation}
        \boldsymbol\xi_{k+1}
        =\mat\Phi_s(\Delta t_k)\boldsymbol\xi_k+\vect\eta_k,
        \qquad
        \vect e^p_{s,k}=\mat H_s\boldsymbol\xi_{k+1}.
        \label{eq:characterization-timing}
    \end{equation}
    The end-of-interval state, not the start prior, affects interval $k$.


    \section{Returned statistics and ordering}
    The result contains one channel estimate for the accelerometer and one for the
    gyroscope. Each channel exposes
    \begin{itemize}
        \item the configured continuous white-noise covariance \(\mat Q_{c,s}\);
        \item interval-average rate covariance;
        \item direct fresh-noise covariance for the next reported increment;
        \item conditional and marginal increment covariance;
        \item deterministic temperature-dependent rate offset and its increment mean;
        \item the persistent-state transition, process covariance, start and end
        priors, state labels/slices, and output measurement matrix;
        \item the increment--state-innovation cross-covariance required by an exact
        discrete augmented filter; and
        \item metadata describing omitted effects and approximations.
    \end{itemize}
    The combined six-dimensional covariance uses the fixed ordering
    \begin{equation}
        \begin{bmatrix}
            \Delta v_x & \Delta v_y & \Delta v_z &
            \Delta\theta_x & \Delta\theta_y & \Delta\theta_z
        \end{bmatrix}^{\mathsf T}.
    \end{equation}
    The current configuration has no accelerometer--gyroscope cross-covariance, so
    the combined covariance is block diagonal. The package reports sensor-domain
    statistics; the navigation filter remains responsible for mapping them through
    its own state-transition and noise-input matrices.

    \begin{table}[htbp]
        \centering
        \caption{Estimator-facing outputs of the stateless characterization.}
        \label{tab:characterization-outputs}
        \begin{tabularx}{\textwidth}{@{}L{1.95in}L{1.42in}Y@{}}
\toprule
\textbf{Output family} & \textbf{Typical symbol} & \textbf{Use} \\
\midrule
Continuous white covariance & $\mat Q_{c,s}$ & Canonical sensor-rate spectral covariance. \\
Interval-average rate covariance & $\mat R^r_{w,s}$ & Fresh rate-noise covariance over the requested sample interval. \\
Persistent-state transition & $\mat\Phi_s(\Delta t)$ & Exact propagation of turn-on, in-run, and flicker state blocks. \\
Persistent process covariance & $\mat Q_s(\Delta t)$ & New uncertainty injected into those persistent blocks. \\
Reset and elapsed-time prior & $\vect\mu_{\xi,s}(t),\mat P_{\xi,s}(t)$ & Configuration-only prior, not a realized bias estimate. \\
Conditional increment covariance & $\mat R^{\Delta}_{s,\mathrm{cond}}$ & Fresh white noise plus optional quantization approximation when bias states are carried. \\
Marginal increment covariance & $\mat P^{\Delta}_{s,\mathrm{marg}}$ & One-interval uncertainty envelope when persistent states are omitted. \\
Known and omitted terms & metadata & Makes temperature context and signal-dependent exclusions explicit. \\
\bottomrule
        \end{tabularx}
    \end{table}


    \section{Persistent statistical state}
    For channel \(s\), the persistent state is represented as three-axis blocks:
    \begin{equation}
        \boldsymbol\xi_s=
        \begin{bmatrix}
            \vect b_{\mathrm{on},s}^{\mathsf T} &
            \vect b_s^{\mathsf T} &
            \vect x_{1,s}^{\mathsf T} &\cdots&
            \vect x_{K,s}^{\mathsf T}
        \end{bmatrix}^{\mathsf T}.
        \label{eq:characterization-state}
    \end{equation}
    Disabled blocks are omitted. The blocks that remain are respectively the
    turn-on bias, the in-run Gauss--Markov or random-walk bias, and the finite-band
    OU components used to approximate flicker bias. The returned
    \software{state\_block\_labels}, \software{state\_block\_slices}, and
    \software{measurement\_matrix} make the state layout explicit:
    \begin{equation}
        \vect e_s=\mat H_s\boldsymbol\xi_s
    \end{equation}
    is the additive persistent rate error seen by the channel. The exact
    discrete-time characterization for one future interval is
    \begin{equation}
        \boldsymbol\xi_{k+1}
        =\mat\Phi_s(\Delta t)\boldsymbol\xi_k+\vect\eta_k,
        \qquad
        \vect\eta_k\sim\Normal(\vect0,\mat Q_s(\Delta t)).
        \label{eq:characterization-discrete-state}
    \end{equation}

    For a stationary Gauss--Markov block with correlation time \(T_b\),
    \begin{align}
        \phi_b&=\exp(-\Delta t/T_b),\\
        \mat\Phi_b(\Delta t)&=\phi_b\mat I,\\
        \mat Q_b(\Delta t)&=(1-\phi_b^2)\mat \Sigma_b.
        \label{eq:characterization-gm}
    \end{align}
    For a random-walk block, \software{bias\_std} is the driving standard
    deviation in sensor units per square root second, and
    \begin{equation}
        \mat\Phi_b(\Delta t)=\mat I,
        \qquad
        \mat Q_b(\Delta t)=\mat \Sigma_{\mathrm{rw}}\Delta t,
        \qquad \mat \Sigma_{\mathrm{rw}}=\diag(\vect\sigma_{\mathrm{rw}}\Had\vect\sigma_{\mathrm{rw}}).
        \label{eq:characterization-rw}
    \end{equation}
    The returned \software{persistent.transition} and
    \software{persistent.process\_covariance} are block-diagonal matrices for the
    full state in \cref{eq:characterization-state}. The convenience properties
    \software{bias\_transition} and \software{bias\_process\_covariance} select the
    three-axis in-run bias block for an estimator that carries only that bias.

    The random-walk density and an initial bias standard deviation have different
    units and are not interchangeable. In the current schema the random-walk state
    starts at a known zero state; reset uncertainty is represented by the separate
    turn-on block. Its covariance therefore grows as
    \begin{equation}
        \mat P_{\mathrm{rw}}(t)=\mat Q_{c,\mathrm{rw}}t,
        \qquad \mat P_{\mathrm{rw}}(0)=\mat0.
    \end{equation}
    A future explicit initial-covariance field may add uncertain random-walk
    initialization without overloading the driving density.


    \section{Start and end priors}
    The characterization exposes the reset prior and both the start and end prior
    \((\boldsymbol\mu_{\xi,0},\mat P_{\xi,0})\) and the prior at the supplied
    elapsed time. The convenience properties
    \software{initial\_bias\_mean}/\software{initial\_bias\_covariance} and
    \software{prior\_bias\_mean}/\software{prior\_bias\_covariance} apply the
    returned output matrix to those two full-state priors. For a time-invariant
    block model,
    \begin{align}
        \boldsymbol\mu_{\xi,k+1}^{-}&=\mat\Phi_s(\Delta t_k)\boldsymbol\mu_{\xi,k},\\
        \mat P_{\xi,k+1}^{-}&=\mat\Phi_s(\Delta t_k)\mat P_{\xi,k}
        \mat\Phi_s(\Delta t_k)^{\mathsf T}+\mat Q_s(\Delta t_k).
        \label{eq:characterization-prior}
    \end{align}
    The configured mean is zero for these stochastic blocks. The reported start and
    end priors are not claims about the realized current bias; a live bias depends on
    reset draws and history, while a filter posterior additionally depends on
    external observations.
    Temperature-dependent deterministic bias is returned separately as
    \software{deterministic\_rate\_offset} and
    \software{deterministic\_increment\_mean}, rather than being confused with a
    random current-bias estimate.


    \section{White-noise and increment covariance}
    Let \(\mat\Lambda_s\) be the diagonal thermal noise multiplier and let
    \(\mat S_s\) be the product of channel and global output scales. The covariance
    of the interval-average rate white noise is
    \begin{equation}
        \mat R^r_{w,s}
        =\frac{\mat\Lambda_s\mat Q_{c,s}\mat\Lambda_s^{\mathsf T}}{\Delta t}.
        \label{eq:characterization-rate-covariance}
    \end{equation}
    The direct covariance of the reported increment is therefore
    \begin{equation}
        \mat R^\Delta_{w,s}
        =\Delta t^2\mat S_s\mat R^r_{w,s}\mat S_s^{\mathsf T}
        =\Delta t\mat S_s\mat\Lambda_s\mat Q_{c,s}
        \mat\Lambda_s^{\mathsf T}\mat S_s^{\mathsf T}.
        \label{eq:characterization-direct-covariance}
    \end{equation}
    This is the covariance for fresh independent white noise on the next
    \(\Delta\vect v\) or \(\Delta\vect\theta\) sample.


    \section{Conditional, joint, and marginal covariance}
    The \software{direct\_increment\_covariance} contains fresh white noise and,
    when requested, the high-resolution quantization approximation. Because the
    state is advanced before use, the start-conditioned covariance also includes
    the persistent innovation generated during the interval:
    \begin{align}
        \mat B_{s,k}&=\Delta t_k\mat S_s\mat H_s,\\
        \mat R^\Delta_{s\mid\xi_k}
        &=\mat R^\Delta_{d,s}+\mat B_{s,k}\mat Q_s\mat B_{s,k}^{\mathsf T},\\
        \mat C_{\Delta\eta,s}&=\mat B_{s,k}\mat Q_s.
        \label{eq:characterization-joint-blocks}
    \end{align}
    The returned \software{increment\_state\_process\_cross\_covariance} is the
    matrix $\mat C_{\Delta\eta,s}$ required when a discrete augmented filter
    injects increment noise and persistent-state process noise separately.
    The corresponding local joint covariance, ordered as increment error followed
    by persistent-state innovation, is
    \begin{equation}
        \mat{\mathcal Q}_{s,k}^{\mathrm{joint}}=
        \begin{bmatrix}
            \mat R^\Delta_{d,s}+\mat B_{s,k}\mat Q_s\mat B_{s,k}^{\mathsf T}
            &\mat B_{s,k}\mat Q_s\\
            \mat Q_s\mat B_{s,k}^{\mathsf T}&\mat Q_s
        \end{bmatrix}.
        \label{eq:characterization-joint}
    \end{equation}

    The \software{marginal\_increment\_covariance} additionally integrates over
    the uncertainty in the persistent prior:
    \begin{equation}
        \mat P^\Delta_{s,\mathrm{marg}}(a_{k+1})
        =\mat R^\Delta_{d,s}
        +\mat B_{s,k}\mat P_{\xi,s,k+1}^{-}\mat B_{s,k}^{\mathsf T}.
        \label{eq:characterization-marginal}
    \end{equation}
    This is useful as a one-step uncertainty envelope when the filter does not
    carry those persistent states. It must not be treated as fresh independent
    noise at every sample: the persistent contribution remains correlated across
    time.

    For two different intervals, the persistent contribution has a nonzero
    cross-covariance. In a time-invariant channel, for $t_j>t_k$,
    \begin{equation}
        \Cov(\Delta\vect y^p_{s,k},\Delta\vect y^p_{s,j})
        =\Delta t_k\Delta t_j\,
        \mat S_s\mat H_s\mat P_{\xi,s}(t_k)
        \mat\Phi_s(t_j-t_k)^\transpose
        \mat H_s^\transpose\mat S_s^\transpose.
        \label{eq:characterization-cross-covariance}
    \end{equation}
    Fresh white noise contributes only to the same interval. Turn-on bias has a
    non-decaying correlation, while Gauss--Markov and flicker components decay
    according to their transitions; random-walk correlation evolves
    nonstationarily. This is why the marginal covariance is an envelope rather
    than a replacement for an augmented colored-error state.

    \begin{tcolorbox}[
        enhanced,breakable,colback=ImuPaleGold,colframe=ImuGold,
        title={Kalman-filter rule},fonttitle=\bfseries]
        For an augmented error-state filter, use the conditional increment covariance
        for fresh sensor noise and propagate the explicit bias state with its returned
        transition and process covariance. Do not also add the full marginal bias
        covariance to every IMU sample. Doing both double-counts the same uncertainty
        and incorrectly turns colored bias into independent white noise.
    \end{tcolorbox}

    \begin{lstlisting}[style=imuPython,caption={Augmented-filter use of sensor-domain statistics.},label={lst:characterization-filter}]
noise = characterize_imu_error_process(
    config=imu_config,
    dt=imu_output.dt,
    process_age_at_interval_start_s=process_age_s,
    temperature_celsius=imu_output.temperature_celsius,
)
estimator.propagate_imu(
    delta_v_body_start=imu_output.delta_v,
    delta_theta_body_start=imu_output.delta_theta,
    dt=imu_output.dt,
    direct_increment_covariance=noise.direct_increment_covariance,
    accel_bias_transition=noise.accelerometer.persistent.bias_transition,
    accel_bias_process_covariance=(
        noise.accelerometer.persistent.bias_process_covariance
    ),
    gyro_bias_transition=noise.gyroscope.persistent.bias_transition,
    gyro_bias_process_covariance=(
        noise.gyroscope.persistent.bias_process_covariance
    ),
    increment_state_process_cross_covariance=(
        noise.increment_state_process_cross_covariance
    ),
)
    \end{lstlisting}


    \section{Continuous-time filter ownership}
    For an error-state filter with local dynamics
    \begin{equation}
        \dot{\delta\vect x}
        =\mat F(t)\delta\vect x+\mat G(t)\vect w(t),
        \qquad
        \E[\vect w(t)\vect w(t')^\transpose]
        =\mat Q_c\delta(t-t'),
        \label{eq:characterization-filter-dynamics}
    \end{equation}
    the complete navigation-state discrete covariance is
    \begin{equation}
        \mat Q_{d,k}
        =\int_0^{\Delta t_k}
        \mat\Phi_x(\Delta t_k,\tau)\mat G(\tau)\mat Q_c
        \mat G(\tau)^\transpose
        \mat\Phi_x(\Delta t_k,\tau)^\transpose\,\dd\tau.
        \label{eq:characterization-filter-covariance}
    \end{equation}
    The matrices in this expression depend on the filter's state ordering, nominal
    attitude, specific force, angular rate, gravity model, and discretization. The
    IMU package therefore owns sensor-domain $\mat Q_c$, persistent-state
    $\mat\Phi$ and $\mat Q$, priors, and increment covariances; the filter owns
    the mapping into its full navigation-state covariance.

    For a stationary Gauss--Markov bias, the corresponding continuous state model
    can be written
    \begin{equation}
        \mat A_b=-\frac{1}{T_b}\mat I,
        \qquad
        \mat Q_{c,b}=\frac{2}{T_b}\mat\Sigma_b,
        \label{eq:characterization-gm-continuous}
    \end{equation}
    while a random walk uses
    \begin{equation}
        \mat A_b=\mat0,
        \qquad
        \mat Q_{c,b}=\mat Q_{c,\mathrm{rw}}.
        \label{eq:characterization-rw-continuous}
    \end{equation}
    The exact local blocks returned by the characterization may be consumed
    directly, or the filter may use these continuous terms in its own Van Loan or
    numerical discretization.


    \section{Deliberate exclusions and filter ownership}
    The function characterizes effects that can be determined from configuration,
    time, and the explicitly supplied temperature. It identifies the following
    effects in metadata rather than inventing a covariance:
    \begin{itemize}
        \item scale factor and radial nonlinearity, which depend on nominal input
        rate;
        \item misalignment covariance, which depends on the current deterministic
        signal;
        \item clipping, which is nonlinear and signal-dependent;
        \item exact quantization behavior, which depends on bin phase and signal
        motion;
        \item realized turn-on and current bias; and
        \item thermal offset and scaling when no temperature is supplied.
    \end{itemize}
    When requested and when a quantization step is configured, the common
    high-resolution approximation is added to the conditional covariance:
    \begin{equation}
        \mat R^r_q\approx\frac{1}{12}\diag\left(\vect\delta_s\Had\vect\delta_s\right),
        \qquad \vect\delta_s=\field{quantization\_step}.
        \label{eq:characterization-quantization}
    \end{equation}
    A scalar configuration value is broadcast to all three axes; a three-axis
    value preserves distinct per-axis quantization steps.
    The complete navigation-state covariance remains filter-owned. For filter state
    transition \(\mat\Phi_x\), noise-input matrix \(\mat G_k\), and continuous
    noise covariance \(\mat Q_c\), the filter computes, for example,
    \begin{equation}
        \mat Q_{d,k}=\int_0^{\Delta t}
        \mat\Phi_x(\Delta t-\tau)\mat G_k\mat Q_c\mat G_k^{\mathsf T}
        \mat\Phi_x(\Delta t-\tau)^{\mathsf T}\,d\tau.
    \end{equation}
    Those matrices depend on the filter's nominal attitude, specific force, angular
    rate, gravity model, and state ordering. The IMU package owns sensor-domain
    \(\mat Q_c\), persistent-state \(\mat\Phi\) and \(\mat Q\), priors, and
    increment covariances; it cannot produce the full navigation covariance from
    configuration alone.

    \begin{table}[htbp]
        \centering
        \caption{Effects that configuration-only characterization deliberately does not
        treat as fresh Gaussian covariance.}
        \label{tab:characterization-limits}
        \begin{tabularx}{\textwidth}{@{}L{1.55in}Y Y@{}}
\toprule
\textbf{Effect} & \textbf{Why configuration and time are insufficient} & \textbf{Treatment} \\
\midrule
Scale and radial nonlinearity & The error depends on the nominal channel rate. & Keep the configured terms as deterministic model parameters or linearize them in the filter when the nominal signal is available. \\
Misalignment & The output covariance depends on the current deterministic signal. & Map a misalignment prior through a signal-dependent Jacobian. \\
Clipping & Saturation is nonlinear and signal-dependent. & Treat the Gaussian covariance as invalid near the configured measurement range. \\
Exact quantization & Bin phase and signal motion determine whether the error is uniform, deterministic, or correlated. & Use only the opt-in $\delta^2/12$ high-resolution approximation. \\
Thermal deterministic offset & A value requires temperature; thermal lag would require additional state. & Evaluate the memoryless model only when Celsius temperature is supplied and keep the offset separate from stochastic bias. \\
Realized current bias & It depends on reset draws, process history, and filter observations. & Obtain it from live state for diagnostics or from the estimator posterior; never infer it from configuration alone. \\
\bottomrule
        \end{tabularx}
    \end{table}


    \section{Implementation architecture}
    The stateful sampler and stateless characterization must describe the same
    process. The intended dependency direction is
    \begin{equation}
        \text{validated configuration}
        \longrightarrow
        \begin{cases}
            \text{shared exact transition helpers},\\
            \text{stateful sampler using those helpers},\\
            \text{stateless characterization using those helpers}.
        \end{cases}
    \end{equation}
    In particular, Gauss--Markov, random-walk, and flicker transition parameters,
    as well as flicker-component construction and calibration, are shared with the
    sampler. The characterization returns newly assembled arrays, so callers can
    consume its matrices without mutating configuration or model-owned state. Its
    thermal evaluation uses the built-in \software{LinearThermalModel}; a custom
    \software{ThermalModel} supplied to \software{ImuModel} is not recoverable from
    \software{ImuConfig} and is therefore outside this config-only contract.

    For efficiency, model construction compiles immutable configuration vectors and
    covariance matrices into cached, read-only NumPy arrays. The sampler reuses
    these arrays for scale factors, bias standard deviations, thermal coefficients,
    quantization steps, and covariance inputs rather than rebuilding them for every
    sample. This cache is an implementation detail: temperature-dependent thermal
    states and timestep-dependent stochastic transitions are still evaluated when
    their inputs change, and the serialized Pydantic configuration remains the
    source of truth.


    \chapter{Verification and validation plan}
    \label{chap:verification}


    \section{Verification philosophy}
    The strongest evidence comes from a hierarchy of tests:
    \begin{enumerate}
        \item exact deterministic identities;
        \item distribution and covariance tests for isolated stochastic terms;
        \item time-correlation and Allan-signature tests;
        \item mixed-effect Monte Carlo cases;
        \item integration tests with a simple dead-reckoning consumer; and
        \item end-to-end closed-loop scenarios.
    \end{enumerate}
    No single Allan plot proves the full model. Conversely, unit tests that only
    exercise formulas at one time step do not establish long-term stochastic
    behavior.


    \section{Deterministic tests}
    A zero-error configuration should satisfy
    \begin{align}
        \Delta\vect v^{\mathrm{meas}}_k&=\Delta\vect v_{b,k},\\
        \Delta\vect\theta^{\mathrm{meas}}_k&=\Delta\vect\theta_{b,k}
        \label{eq:zero-error-identity}
    \end{align}
    within floating-point tolerance when output scales are unity and clipping and
    quantization are disabled.

    Frame tests should include
    \begin{itemize}
        \item identity attitude and a known world velocity step;
        \item a \SI{90}{\degree} body rotation that maps one world axis into another;
        \item invariance of relative rotation under a fixed world-frame rotation;
        \item positive and negative small-angle increments; and
        \item rotations near \(\pi\), where logarithm implementations are most
        numerically delicate.
    \end{itemize}

    For each deterministic error stage, use basis-vector and mixed-axis inputs to
    verify scale, nonlinearity, thermal bias, thermal scale, and misalignment
    independently. Finite differences can check \cref{eq:deterministic-jacobian,eq:misalignment-jacobian,eq:thermal-jacobian} away from clip and quantization
    boundaries.


    \section{White-noise tests}
    For \(N\) independent rate samples at fixed \(\Delta t\), estimate
    \begin{equation}
        \widehat{\mat Q}_c
        =\Delta t\,\widehat{\Cov}(\vect n_k).
        \label{eq:empirical-qc}
    \end{equation}
    The result should converge to the configured \(\mat Q_c\). Run the same test at
    multiple \(\Delta t\) values: rate variance should scale as \(1/\Delta t\), while
    increment variance should scale as \(\Delta t\).

    For a scalar normal variance estimate from \(N\) independent samples,
    \begin{equation}
        \frac{(N-1)\widehat\sigma^2}{\sigma^2}\sim\chi^2_{N-1},
        \label{eq:variance-chi-square}
    \end{equation}
    which supplies a principled acceptance interval instead of an arbitrary percent
    error. Multivariate covariance tests can compare eigenvalues, correlations, and
    Mahalanobis statistics.


    \section{Gauss--Markov and random-walk tests}
    A stationary Gauss--Markov fixture should verify
    \begin{align}
        \widehat{\Var}(b_i)&\approx\sigma_b^2,\\
        \widehat\rho(\ell)&\approx e^{-\ell\Delta t/T},
        \label{eq:gm-test-targets}
    \end{align}
    including at least one variable-step schedule. A random-walk fixture should
    verify that increment variance is
    \begin{equation}
        \Var(b_{k+1}-b_k)=\sigma_{\mathrm{rw}}^2\Delta t_k
    \end{equation}
    and that ensemble variance grows linearly with elapsed time.


    \section{Flicker and Allan tests}
    The compact regression fixtures under \software{tests/profiles/test/} should
    compare
    \begin{itemize}
        \item white-noise level at short averaging times;
        \item Gauss--Markov turnover location and asymptotic slopes;
        \item random-walk long-time slope; and
        \item finite-band flicker level near \(\tau_c\).
    \end{itemize}
    Tests should compare bands or fitted slopes rather than demand pointwise equality
    to analytic curves. Multiple seeds reduce the chance that one favorable record
    masks an incorrect implementation.


    \section{Reset and reproducibility tests}
    With a fixed seed and identical sample history, every output should reproduce
    bit-for-bit when the numerical backend permits. Across many distinct seeds,
    reset-level samples should satisfy their configured covariance. A reset should
    also clear prior truth history so the next call returns zero increments rather
    than bridging across runs.


    \section{Output-electronics tests}
    Clipping tests should cover below-range, exact-boundary, and over-range values on
    both signs. Quantization tests should cover exact bins, values just inside each
    side of a threshold, and half-step ties. A mixed test should confirm that
    clipping precedes quantization and that both precede multiplication by
    \(\Delta t\) and output scales.


    \section{Stateless characterization tests}
    The characterization path should be tested independently of random trajectories.
    For fixed configuration and independent inputs, repeated calls must return
    identical matrices and must leave the configuration unchanged. Exact identities
    should cover
    \begin{itemize}
        \item Gauss--Markov \(\phi\) and \(Q(\Delta t)\);
        \item random-walk transition and elapsed-time covariance growth;
        \item the reset-to-elapsed prior covariance equation;
        \item direct white increment covariance and output scaling;
        \item conditional versus marginal covariance separation; and
        \item positive-semidefiniteness and the documented six-dimensional ordering.
    \end{itemize}
    Quantization tests should verify that the \(\delta^2/12\) term is opt-in and
    explicitly identified in metadata. A filter integration test should verify that
    conditional covariance plus explicit bias propagation does not also consume the
    marginal persistent covariance.


    \section{Verification matrix}
    \begin{longtable}{@{}L{1.45in}L{2.05in}L{2.75in}@{}}
        \caption{Recommended verification matrix and exit evidence.}
        \label{tab:verification-matrix}\\
        \toprule
        \textbf{Capability} & \textbf{Primary test} & \textbf{Exit evidence} \\
        \midrule
        \endfirsthead
        \toprule
        \textbf{Capability} & \textbf{Primary test} & \textbf{Exit evidence} \\
        \midrule
        \endhead
        Truth-frame contract & Analytic rotations and velocity steps & Correct axes, signs, units, and world-rotation invariance. \\
        Zero-error pass-through & Deterministic identity fixture & Output equals \cref{eq:ideal-dv,eq:ideal-dtheta}. \\
        Scale and nonlinearity & Basis and mixed-axis sweeps & Formula parity and finite-difference Jacobian agreement. \\
        Thermal behavior & Multi-temperature sweep and \software{None} case & Correct bias, scale, noise multiplier, reference temperature, and identity state. \\
        Misalignment & Fixed seeded rotation & Norm preservation and expected first-order cross-axis coupling. \\
        White noise & Multi-$\Delta t$ covariance test & $\Delta t\widehat{\Cov}(n)\approx Q_c$ and correct thermal covariance. \\
        Turn-on bias & Reset ensemble & Configured mean/covariance; constant inside each run. \\
        Gauss--Markov & Long stationary run & Stationary variance and exponential autocorrelation. \\
        Random walk & Multi-horizon ensemble & Increment and elapsed-time variance laws. \\
        Flicker approximation & Allan-band regression & Target near $\tau_c$ and roll-off outside configured band. \\
        Stateless characterization & Pure-function and matrix-identity tests & Exact transitions, priors, conditional/marginal separation, and six-dimensional ordering. \\
        Clipping & Boundary table & Exact componentwise saturation and disabled identity behavior. \\
        Quantization & Bin and tie table & Exact configured rounding rule and error bound. \\
        Output scaling & Known deterministic rate & Correct channel/global scales and integration by $\Delta t$. \\
        Reset reproducibility & Same-seed replay & Identical reset state and output sequence. \\
        Protocol substitution & Alternate dummy implementation & Consumer depends on \software{ImuModelProtocol}, not concrete class internals. \\
        Packaging & Build/install smoke test & Wheel/sdist install, import, examples, and metadata succeed. \\
        \bottomrule
    \end{longtable}


    \section{Verification criteria}
    The implementation and its examples should satisfy the following criteria:
    \begin{itemize}
        \item all deterministic contracts have exact or tolerance-bounded unit tests;
        \item statistical tests use documented confidence or Monte Carlo tolerances;
        \item at least one white, Gauss--Markov, random-walk, and flicker fixture has a
        saved analytic reference;
        \item public inputs, outputs, units, frame directions, and reset behavior match
        this specification;
        \item examples distinguish notional profiles from authoritative specifications;
        \item a small downstream integration demonstrates that a consumer can use
        \field{delta\_v}, \field{delta\_theta}, and \field{dt} without reading truth.
    \end{itemize}


    \chapter{Assumptions, limitations, and roadmap}
    \label{chap:limitations}


    \section{Current modeling assumptions}
    The default model assumes
    \begin{itemize}
        \item one timestamp and one truth velocity/orientation snapshot per sample;
        \item accelerometer and gyro error channels are independent unless covariance
        is configured within a channel;
        \item run-level misalignment is a proper rotation and remains fixed;
        \item deterministic nonlinearity is radial and quadratic;
        \item the default thermal response is instantaneous and linear in temperature;
        \item white-noise parameters represent continuous-time density or covariance;
        \item Gauss--Markov states use exact first-order transitions;
        \item flicker behavior is finite-band and approximated by a finite component
        bank; and
        \item clipping and quantization occur in the rate domain.
    \end{itemize}


    \section{Not represented by the default model}
    The following effects require preprocessing, an alternate implementation, or a
    future extension:
    \begin{itemize}
        \item gravity and Earth-rate terms;
        \item lever-arm acceleration, centripetal acceleration, and mounting flex;
        \item finite sensor bandwidth, anti-aliasing filters, group delay, and sample
        aperture;
        \item coning and sculling corrections from subinterval motion;
        \item gyro acceleration sensitivity and vibration rectification;
        \item hysteresis, aging, warm-up transients, shock recovery, and radiation;
        \item cross-channel accelerometer--gyro covariance and common environmental
        latent processes;
        \item temperature dynamics, self-heating, and spatial gradients;
        \item calibration estimation, online bias estimation, or navigation filtering;
        and
        \item vendor-certified profiles.
    \end{itemize}


    \section{Candidate roadmap}
    A technically coherent roadmap is
    \begin{enumerate}
        \item \textbf{Core hardening}: freeze frame, unit, reset, and stochastic-index
        semantics; complete analytic regression fixtures.
        \item \textbf{Profile tooling}: add schema versions, provenance fields, unit
        conversion reports, and evidence-quality labels.
        \item \textbf{Dynamic response}: add optional linear filters and explicit
        sample aperture while preserving continuous-time noise semantics.
        \item \textbf{High-dynamic inertial increments}: add substep truth input and
        coning/sculling-aware models.
        \item \textbf{Environmental coupling}: add thermal lag, vibration sensitivity,
        and shared cross-channel latent processes.
        \item \textbf{Calibration and fitting}: estimate parameters from static records,
        temperature sweeps, turntable tests, and Allan data with uncertainty reports.
        \item \textbf{Integration evidence}: provide estimator and closed-loop examples
        that quantify degradation from ideal truth to progressively richer error
        profiles.
    \end{enumerate}

    \begin{engineeringbox}
        Extensions should add new, explicitly named operators rather than overloading an
        existing scalar with a different physical meaning. Backward-compatible clarity is
        more valuable than compressing every hardware behavior into one configuration
        object.
    \end{engineeringbox}


    \appendix


    \chapter{Mathematical derivations}
    \label{app:derivations}


    \section{Exact Gauss--Markov discretization}
    Starting from the scalar form of \cref{eq:ou-sde},
    \begin{equation}
        \dd b(t)=-\frac{1}{T}b(t)\dd t+q\dd W(t),
        \qquad q=\sqrt{\frac{2\sigma_b^2}{T}},
    \end{equation}
    the solution over \(\Delta t\) is
    \begin{equation}
        b(t+\Delta t)
        =e^{-\Delta t/T}b(t)
        +q\int_t^{t+\Delta t}
        e^{-(t+\Delta t-\tau)/T}\dd W(\tau).
        \label{eq:ou-solution}
    \end{equation}
    The stochastic integral is Gaussian with variance
    \begin{align}
        Q_d
        &=q^2\int_t^{t+\Delta t}
        e^{-2(t+\Delta t-\tau)/T}\dd\tau\\
        &=\frac{2\sigma_b^2}{T}\cdot\frac{T}{2}
        \left(1-e^{-2\Delta t/T}\right)\\
        &=\sigma_b^2(1-\phi^2),
        \qquad \phi=e^{-\Delta t/T},
    \end{align}
    which yields \cref{eq:gm-discrete}. If
    \(\Var(b_k)=\sigma_b^2\), then
    \begin{equation}
        \Var(b_{k+1})
        =\phi^2\sigma_b^2+\sigma_b^2(1-\phi^2)
        =\sigma_b^2,
    \end{equation}
    proving stationarity under the exact update.


    \section{White-noise density to increment variance}
    From \cref{eq:white-discrete-cov}, a scalar rate sample has variance
    \(N^2/\Delta t\). Its increment contribution is
    \begin{equation}
        \Delta y_k=\Delta t\,n_k,
    \end{equation}
    so
    \begin{equation}
        \Var(\Delta y_k)
        =\Delta t^2\frac{N^2}{\Delta t}
        =N^2\Delta t.
    \end{equation}
    Over \(M\) independent intervals of equal duration, the integrated noise
    variance is
    \begin{equation}
        \Var\!\left(\sum_{k=0}^{M-1}\Delta y_k\right)
        =MN^2\Delta t=N^2t,
        \qquad t=M\Delta t,
    \end{equation}
    which recovers continuous random-walk growth.


    \section{Gauss--Markov Allan variance}
    For a stationary scalar process with autocovariance
    \begin{equation}
        R_b(\tau)=\sigma_b^2e^{-\lvert\tau\rvert/T},
    \end{equation}
    the variance of a \(\tau\)-length average is
    \begin{align}
        \Var(\bar b)
        &=\frac{1}{\tau^2}\int_0^\tau\!\int_0^\tau
        R_b(t-s)\dd t\dd s\\
        &=\frac{2\sigma_b^2}{\tau^2}
        \left[T\tau-T^2(1-e^{-\tau/T})\right].
        \label{eq:gm-mean-var}
    \end{align}
    The covariance of adjacent averages is
    \begin{equation}
        \Cov(\bar b_j,\bar b_{j+1})
        =\frac{\sigma_b^2T^2}{\tau^2}(1-e^{-\tau/T})^2.
        \label{eq:gm-adjacent-cov}
    \end{equation}
    Since Allan variance is half the variance of the difference between adjacent
    averages,
    \begin{align}
        \sigma_A^2(\tau)
        &=\Var(\bar b)-\Cov(\bar b_j,\bar b_{j+1})\\
        &=\frac{
            2\sigma_b^2[T\tau-T^2(1-e^{-\tau/T})]
            -\sigma_b^2T^2(1-e^{-\tau/T})^2
        }{\tau^2},
    \end{align}
    which is equivalent to \cref{eq:allan-gm}.


    \section{Small-angle misalignment}
    Using the first terms of \cref{eq:rodrigues},
    \begin{equation}
        \mat M=\Exp(\skewmat{\vect\epsilon})
        =\mat I+\skewmat{\vect\epsilon}+\mathcal O(\lVert\vect\epsilon\rVert^2).
    \end{equation}
    Therefore
    \begin{equation}
        \mat M\vect d
        =\vect d+\skewmat{\vect\epsilon}\vect d+\mathcal O(\epsilon^2)
        =\vect d-\skewmat{\vect d}\vect\epsilon+\mathcal O(\epsilon^2).
    \end{equation}
    The off-axis error magnitude obeys
    \begin{equation}
        \lVert\vect\epsilon\times\vect d\rVert
        =\lVert\vect\epsilon\rVert\lVert\vect d\rVert\sin\alpha,
    \end{equation}
    where \(\alpha\) is the angle between misalignment and signal vectors. A
    single-axis calibration cannot fully reveal a general three-axis misalignment.


    \section{Random-walk Allan variance}
    Let \(b(t)=\sigma_{\mathrm{rw}}W(t)\). For adjacent interval averages of length
    \(\tau\), direct integration of Brownian covariance
    \(\E[W(t)W(s)]=\min(t,s)\) gives
    \begin{equation}
        \frac12\Var(\bar b_{j+1}-\bar b_j)
        =\frac{\sigma_{\mathrm{rw}}^2\tau}{3},
    \end{equation}
    which yields \cref{eq:allan-rw}.


    \chapter{Configuration review checklist}
    \label{app:checklist}

    Before accepting a profile or scenario, review the following items.


    \section*{Frames and sampling}
    \begin{itemize}
        \item Is the supplied orientation definitely body-to-world?
        \item Is the velocity input world-frame and gravity-excluded under the intended
        truth adapter?
        \item Are timestamps strictly increasing and in seconds?
        \item Is the sample interval small enough for endpoint differencing in the
        intended maneuvers?
        \item Does the consumer understand that both outputs use body axes at interval
        start?
    \end{itemize}


    \section*{Units and provenance}
    \begin{itemize}
        \item Has every value been converted to canonical SI exactly once?
        \item Are original units and conversion equations retained?
        \item Are typical, maximum, RMS, and one-sigma values distinguished?
        \item Are noise density, Allan coefficient, and integrated random-walk metrics
        distinguished?
        \item Is every estimated value labeled as estimated rather than official?
    \end{itemize}


    \section*{Stochastic semantics}
    \begin{itemize}
        \item Is turn-on bias separated from in-run drift?
        \item If \field{bias\_correlation\_time} is absent, is the random-walk driving
        interpretation intentional?
        \item If a covariance is supplied, is it symmetric positive semidefinite and in
        continuous-time units?
        \item Does the flicker correlation-time band cover the averaging times of
        interest?
        \item Are seeds, reset points, and Monte Carlo trial identifiers recorded?
    \end{itemize}


    \section*{Output behavior}
    \begin{itemize}
        \item Are clipping limits rate-domain limits in canonical units?
        \item Is the quantization step a true output LSB rather than a marketing
        resolution metric?
        \item Are channel and global output scales both intentional?
        \item Is the configured rounding rule covered by tests?
        \item Are temperature-unavailable samples intentionally represented by
        \software{None}?
    \end{itemize}


    \chapter{Requirements traceability summary}
    \label{app:traceability}

    \begin{longtable}{@{}L{0.75in}L{2.45in}L{3.05in}@{}}
        \caption{High-level specification requirements.}
        \label{tab:requirements}\\
        \toprule
        \textbf{ID} & \textbf{Requirement} & \textbf{Primary evidence} \\
        \midrule
        \endfirsthead
        \toprule
        \textbf{ID} & \textbf{Requirement} & \textbf{Primary evidence} \\
        \midrule
        \endhead
        IMU-001 & The model shall accept body-to-world orientation and gravity-excluded world velocity. & \Cref{chap:truth}; frame unit tests. \\
        IMU-002 & The first sample after construction or reset shall initialize history and return zero increments. & \Cref{chap:algorithm}; lifecycle tests. \\
        IMU-003 & Later timestamps shall be strictly increasing. & Input validation tests. \\
        IMU-004 & Ideal increments shall be expressed in body axes at interval start. & \Cref{eq:ideal-dv,eq:ideal-dtheta}; analytic frame fixtures. \\
        IMU-005 & Thermal coefficients shall be bound to a channel model and evaluated from optional temperature. & \Cref{eq:thermal-bias,eq:thermal-identity}; thermal sweep tests. \\
        IMU-006 & Deterministic channel distortion shall include configured scale, radial nonlinearity, thermal scale, and thermal bias. & \Cref{eq:deterministic-map}; formula parity tests. \\
        IMU-007 & Run-level misalignment shall be sampled at reset and held fixed. & \Cref{eq:misalignment}; same-run/reset ensemble tests. \\
        IMU-008 & White rate-noise variance shall scale inversely with sample interval. & \Cref{eq:white-discrete-cov}; multi-rate covariance tests. \\
        IMU-009 & A configured continuous covariance shall replace the independent-axis density path. & \Cref{eq:general-white}; correlated covariance tests. \\
        IMU-010 & Stationary bias with correlation time shall use the exact Gauss--Markov transition. & \Cref{eq:gm-discrete}; stationarity/autocorrelation tests. \\
        IMU-011 & Bias without correlation time shall use the random-walk transition. & \Cref{eq:random-walk}; ensemble variance-growth tests. \\
        IMU-012 & Flicker-like bias shall use a finite logarithmic bank normalized to an Allan target. & \Cref{eq:flicker-sum,eq:flicker-normalization}; Allan-band tests. \\
        IMU-013 & Clipping and quantization shall operate on rate before integration. & \Cref{eq:clip,eq:quantization,eq:unified-model}; boundary tests. \\
        IMU-014 & Output shall include measured increments and interval metadata. & Public contract and integration smoke tests. \\
        IMU-015 & Alternate implementations shall be usable through \software{ImuModelProtocol}. & Protocol conformance test with a dummy model. \\
        IMU-016 & Hardware-oriented profiles shall identify provenance and shall not imply certification. & Profile schema validation and documentation review. \\
        \bottomrule
    \end{longtable}


    \chapter{Compact equation reference}
    \label{app:quick-reference}

    \begin{tcolorbox}[
        enhanced,breakable,colback=ImuPale,colframe=ImuBlue,
        title={Kinematics},fonttitle=\bfseries]
        \begin{align*}
            \Delta t_k&=t_{k+1}-t_k>0,\\
            \Delta\vect v_{b,k}
            &=\Rwb(t_k)^\transpose[\vect v^w(t_{k+1})-\vect v^w(t_k)],\\
            \Delta\vect\theta_{b,k}
            &=\Log[\Rwb(t_k)^\transpose\Rwb(t_{k+1})]^\vee,\\
            \vect u_{a,k}&=\Delta\vect v_{b,k}/\Delta t_k,
            \qquad
            \vect u_{g,k}=\Delta\vect\theta_{b,k}/\Delta t_k.
        \end{align*}
    \end{tcolorbox}

    \begin{tcolorbox}[
        enhanced,breakable,colback=ImuPaleTeal,colframe=ImuTeal,
        title={Deterministic and output model},fonttitle=\bfseries]
        \begin{align*}
            \vect d_s
            &=\vect u_s\Had(\vect1+\vect s_s+\vect\kappa_{T,s})
            \Had(\vect1+n_s\lVert\vect u_s\rVert^2\vect1)+\vect\beta_{T,s},\\
            \mat M_s&=\Exp(\skewmat{\vect\epsilon_s}),\\
            \vect r_{s,k}
            &=\mat M_s\vect d_s+\vect b_{\mathrm{on},s}
            +\vect b_{s,k}+\vect f_{s,k}+\vect n_{s,k},\\
            \Delta\vect y^{\mathrm{meas}}_{s,k}
            &=\Delta t_k\,o_sO_s\,
            \mathcal Q_{\vect\delta_s}(\mathcal C_{L_s}(\vect r_{s,k})).
        \end{align*}
    \end{tcolorbox}

    \begin{tcolorbox}[
        enhanced,breakable,colback=ImuPaleGold,colframe=ImuGold,
        title={Stochastic updates},fonttitle=\bfseries]
        \begin{align*}
            \vect n_{s,k}&\sim\Normal\!\left(
                                           \vect0,\frac{\mat\Lambda_{s,k}\mat Q_{c,s}
                    \mat\Lambda_{s,k}^\transpose}{\Delta t_k}\right),\\
            \vect b_{k+1}&=e^{-\Delta t/T}\vect b_k
            +\Normal\!\left(\vect0,
                          \sigma_b^2(1-e^{-2\Delta t/T})\mat I\right),\\
            \vect b_{k+1}&=\vect b_k
            +\Normal\!\left(\vect0,
                          \sigma_{\mathrm{rw}}^2\Delta t\mat I\right)
            \quad\text{(random walk)},\\
            \vect f(t)&=\sum_{i=1}^{K}\vect x_i(t),
            \qquad T_i\in\operatorname{geomspace}(T_{\min},T_{\max},K).
        \end{align*}
    \end{tcolorbox}

    \begin{tcolorbox}[
        enhanced,breakable,colback=white,colframe=ImuOrange,
        title={Allan references},fonttitle=\bfseries]
        \begin{align*}
            \sigma_{A,\mathrm{white}}(\tau)&=N/\sqrt\tau,\\
            \sigma_{A,\mathrm{rw}}(\tau)&=\sigma_{\mathrm{rw}}\sqrt{\tau/3},\\
            \sigma_{A,\mathrm{GM}}^2(\tau)
            &=\sigma_b^2\left[
                            \frac{2T}{\tau}-\left(\frac{T}{\tau}\right)^2
                (3-4e^{-\tau/T}+e^{-2\tau/T})
            \right].
        \end{align*}
    \end{tcolorbox}

    \printbibliography[heading=bibintoc,title={References}]

\end{document}
